\documentclass[11pt,letterpaper]{article}
\usepackage[T1]{fontenc}
\usepackage{lmodern}
\usepackage[margin=1in]{geometry}
\usepackage{amsmath,amssymb,amsthm,mathtools}
\usepackage{microtype,booktabs,array,longtable,enumitem}
\usepackage[numbers,sort&compress]{natbib}
\usepackage[colorlinks=true,linkcolor=black,citecolor=black,urlcolor=blue,breaklinks=true]{hyperref}
\usepackage{fancyhdr}
\pagestyle{fancy}\fancyhf{}\fancyhead[L]{\small Finite information cuts: precision, range, and depth}\fancyhead[R]{\small Research synthesis}\fancyfoot[C]{\thepage}
\setlength{\headheight}{14pt}
\setlength{\parskip}{3pt}
\clubpenalty=10000
\widowpenalty=10000
\displaywidowpenalty=10000
\setlist{nosep,leftmargin=*}
\newtheorem{theorem}{Theorem}[section]
\newtheorem{proposition}[theorem]{Proposition}
\newtheorem{lemma}[theorem]{Lemma}
\newtheorem{corollary}[theorem]{Corollary}
\theoremstyle{definition}\newtheorem{definition}[theorem]{Definition}
\theoremstyle{remark}\newtheorem{remark}[theorem]{Remark}
\newcommand{\TV}{\operatorname{TV}}
\newcommand{\E}{\mathbb{E}}
\newcommand{\Prob}{\mathbb{P}}
\newcommand{\Var}{\operatorname{Var}}
\newcommand{\R}{\mathbb{R}}
\newcommand{\ind}{\mathbf{1}}
\newcommand{\floor}[1]{\left\lfloor #1\right\rfloor}
\newcommand{\ceil}[1]{\left\lceil #1\right\rceil}
\newcommand{\betaz}{\beta_0}
\newcommand{\Risk}{\mathcal{R}}
\newcommand{\Dclass}{\mathcal{D}_2(\beta)}
\newcommand{\Lip}{\operatorname{Lip}}
\newcommand{\supp}{\operatorname{supp}}
\newcommand{\bl}{\operatorname{bitlength}}
\newcommand{\unit}[1]{\textsf{#1}}
\newcommand{\source}[1]{\textnormal{[dossier: \unit{#1}]}}
\hypersetup{pdftitle={Finite Information Cuts in Markov Prediction: Precision, Range, and Depth},pdfsubject={AI-assisted mathematical research synthesis},pdfauthor={},pdfkeywords={Markov prediction, causal attention, finite precision, quantization, minimax risk}}
\title{\vspace{-1.4em}Finite Information Cuts in Markov Prediction:\\Precision, Range, and Depth}
\author{}
\date{Research synthesis revision 4\quad 2 November 2026}
\begin{document}
\maketitle\vspace{-1.4em}
\begin{abstract}
We study prediction of an unknown homogeneous second-order source at its actual terminal context through a complete finite information cut. The decoder retains quantized summaries and bounded endpoint information, but no history reread or clean statistic. For a fixed-dimensional, fixed-locality, one-summary-block positive-weight family, the minimax total-variation risk is of order
\[
 \max\left\{M^{-1/2},Q^{-1},
 \frac{1+\ln(1+\kappa Q/(M\phi_K))}{\kappa Q}\right\},\qquad
 \phi_K=
 \begin{cases}
 \max\{\sqrt\beta,M^{-1}\},&K=2,\\
 \max\{\beta,M^{-1}\},&K\ge3.
 \end{cases}
\]
Here $Q$ is final-grid resolution, $\kappa$ bounds the weight ratio, and $\beta$ indexes an all-state two-step minorization condition with a hidden stationary certificate. The result holds under stationary and arbitrary initialization. A separately specified finite evaluator of fixed-depth bounded-score causal affine/ReLU encoders has the corresponding profile at $\kappa=1$. At $\beta=1/M$ and $Q\asymp\sqrt M$, binary sources require only constant weight ratio, versus $\Theta(\log M)$ for each fixed larger alphabet. The arbitrary-start binary proof uses reciprocal return identities, first-hit densities, weighted potentials and absorbing-run generating functions, not endpoint mixing. Finally, with uniformly bounded primitive budgets, a zero-logit dyadic compander attains sampling-order risk at $Q\asymp\sqrt M$ using $O(\log\log M)$ actual primitive depth. A rescaled late-window lower proves the matching necessary depth order for $K\ge3$ and $\beta=0$. Working precision, parameter storage and finite evaluation are all charged. These are scoped information-cut results, not universal neural-network or total-memory optimality claims; the manuscript identifies its inherited statistical inputs and complete modular proof dossier.
\end{abstract}
\noindent\textbf{Status and scope.} This is an AI-assisted research synthesis with independent LLM mathematical audits. It is not externally peer reviewed, formally verified, or a submission-ready standalone proof: the full statistical and implementation proof chains are in the accompanying dossier. No authorship or institutional affiliation is assigned. The scientific scope of this revision is frozen at 06:56 UTC on 2 November 2026.

\section{Introduction}
A normalized attention vector can retain a conditional empirical distribution exactly in real arithmetic while losing statistically important information after coordinate rounding. The difficulty is not simply the number of bits in a prediction. A rare context may contribute only a small fraction of the global attention mass. Dividing by that fraction after a finite information cut amplifies the rounding error; more post-cut computation cannot reconstruct distinctions that the complete message no longer contains.

This manuscript studies that phenomenon in an explicit stochastic experiment. The source table is unknown, transitions may vanish, initialization need not be stationary, and the loss is measured at the actual final pair of the same trajectory used to construct the summary. These choices exclude an independent-query simplification. They also make the occupation count of the terminal context a central random quantity. A mixing promise reduces the numerical penalty, but neither the encoder nor the constructive decoder needs the unknown stationary law or a runtime mixing-parameter estimate.

The first resource axis is grid resolution. An original finite normalized-statistic interface with $9K+1$ retained scalars has a complete all-resolution risk profile. For larger alphabets that profile remains valid after optimizing the encoder within a specified fixed-depth causal family evaluated by a certified finite arithmetic procedure. The latter assertion requires more than an analogy with bounded attention: it requires a causal influence estimate, a complete-message concentration and covering argument, and an actual finite-arithmetic upper construction belonging to the family.

The second axis is attention-weight range. The representation can also be changed. A complete $K^3$ triple histogram, with a positive integer amplification of the queried pair before quantization, attains sampling-order risk at square-root grid resolution. A rescaled rare-window proof and a weight-envelope-uniform IID floor give a complete joint weight/grid frontier in the fixed-dimensional, fixed-locality, one-global-summary family. Increasing the weight ratio cannot eliminate the separate $1/Q$ floor. Statistical-order prediction requires both $Q\gtrsim\sqrt M$ and a mixing-dependent lower bound on the product $\kappa Q$.

The alphabet changes that mixing dependence. A binary active Markov segment yields a lower with cutoff $\max\{\sqrt\beta,M^{-1}\}$. A stationary weighted-return analysis first matches it; a path-sensitive first-hit and occupation argument then extends the same paid decoder to every initialization, including off-support starts. At $\beta=1/M$, binary sources need only constant weight ratio in a square-root grid band, whereas each fixed $K\ge3$ needs order $\log M$. At zero mixing promise both require logarithmic ratio. An exact-label histogram inclusion establishes the binary refinement separately for the bounded causal finite family.

The third axis is actual bounded-primitive depth. A fixed-width dyadic compander with zero attention logits achieves statistical-order risk at square-root grid resolution with $O(\log\log M)$ primitive stages, uniformly in the source's initialization and mixing promise. A depth-independent rescaled late-window tube proves necessity of this order for $K\ge3$ at $\beta=0$. Both statements use a genuinely uniform growing-depth arithmetic certificate; neither hides increasing complexity inside a per-layer graph or a fixed architecture constant.

The development proceeds from the experiment and bounded-grid theorem (Sections~\ref{sec:experiment}--\ref{sec:main}), through complete-message, causal, arithmetic and statistical proofs, to the joint weight/grid frontier (Section~\ref{sec:joint}). The binary stationary argument is a useful intermediate step (Section~\ref{sec:binary}); Sections~\ref{sec:initialization}--\ref{sec:binarycausal} remove its initialization restriction. Section~\ref{sec:depth} treats the depth axis. The resource ledger and proof map keep the different implementations and inherited inputs explicit. None of these results is a constant-total-precision implementation or a computational advantage over direct counting.

\subsection{Position relative to existing constructions and lower bounds}
The counting representation is established prior art. Behrens, Biggio and Zdeborov\'a give a weighted inventory with equal-token and unequal-token weights $e$ and $1$, selected mass $ek/[L+(e-1)k]$, and an inversion formula, in an infinite-precision model \citep[Appendix B.3]{behrens2025}. Yehudai et al. study full histogram banks and query-count inversion \citep[Sections 4.1--4.2]{yehudai2026}. Grouping a local tuple as one categorical feature and normalizing a selected sub-block is not a new counting method. The questions here concern the complete post-summary grid, terminal-context occupation, and a lower bound uniform in a growing weight ratio.

Conditional empirical Markov prediction by shallow attention also predates this synthesis. Rajaraman et al. construct constant-depth conditional $k$-gram representations and state logarithmic-bit approximations for their attention-only constructions; their separate layer-normalized construction does not carry the same finite-error analysis \citep[Theorems 4.4 and 5.1; Remarks 4.2 and 5.2]{rajaraman2024}. Real-arithmetic representability and a paid all-operations implementation are therefore distinct claims, not competing descriptions of one theorem.

Bounded-score dilution is likewise a known mechanism. Hahn's sensitivity analysis must be read with the 2021 correction retaining attention-weight variation \citep{hahn2020,hahn2021supp}. Chiang and Cholak show why small continuous sensitivity is not itself information loss: exact recognition can survive shrinking margins, and zero-epsilon normalization or length-scaled logits can evade bounded-Lipschitz assumptions \citep{chiang2022}. Our decoder is unrestricted after the cut, so its impossibility result must follow from the complete finite message, not from a Lipschitz readout hypothesis.

The statistical tools have close predecessors. Han, Jana and Wu use waiting-state sources, a common prior and harmonic averaging in stationary full-history KL prediction; their spectral-gap results also show a mixing-dependent disappearance of a prediction logarithm \citep{han2022}. Those are not same-rate TV results for a quantized interface. Xu and Raginsky give finite-message Bayes lower bounds and the half-logarithmic communication threshold for Bernoulli estimation \citep[Theorem 3 and Example 6]{xu2016}. A generic count of $Q^D$ labels, however, is too weak for our $D$-coordinate bank. The specific obligation is to prove that, with high probability, its \emph{complete} message occupies only $O(1+Qj/M)$ labels on a rare window of length $j$.

Nonlinear scalar quantization is also established. Adler, Tang and Polyanskiy study probability companders, including power transforms $x^s$ and $s=1/2$, followed by inverse reconstruction and normalization under KL distortion \citep[Sections I--III, Theorem 5]{adler2022}. Our depth construction does not introduce companding or square-root behavior near zero; it supplies an explicit bounded-primitive dyadic geometry, one-sided arithmetic and actual-terminal conditional-row analysis. Ou, Schenkel and B\"olcskei construct a depth--weight-precision conversion for ReLU networks \citep[Proposition 4.1]{ou2024}. That coefficient-precision experiment differs from a quantized observed-summary cut with all internal bits paid. The cited statements and immediate explanations were inspected; their full appendix proofs are not claimed to have been audited here.

Accordingly, the contribution asserted here is a source- and interface-specific synthesis of established mechanisms with explicit additional obligations: complete-message accounting, causal depth, dependence-safe terminal occupation, finite arithmetic, and quantitative growing-range constants. The literature comparison is targeted, not a certificate of novelty or an exhaustive survey.

\section{Experiment and information interfaces}\label{sec:experiment}
Fix an alphabet $\mathcal A$ of size $K\ge2$, $M\ge64$, and a dyadic resolution $Q=2^h$ with integer $h\ge4$. Observe $X_1,\ldots,X_{M+2}$ from an unknown homogeneous second-order table $\theta_{ab}(c)$. Set
\[
 Y_j=(X_{j+1},X_{j+2}),\quad 0\le j\le M,
 \qquad P((a,b),(b,c))=\theta_{ab}(c).
\]
The target is $\theta_{Y_M}$, the next-symbol law at the observed terminal pair. A predictor $\widehat p$ has risk
\begin{equation}\label{eq:risk}
 \E_{\theta,\mu}\TV(\widehat p,\theta_{Y_M}),
 \qquad \TV(p,q)=\frac12\sum_{c\in\mathcal A}|p_c-q_c|,
\end{equation}
where $\mu$ is the initial pair law. All expectations are unconditional on the actual trajectory; $Y_M$ is not an independent query.

\begin{definition}[All-state certificate]
The class $\Dclass$, $0\le\beta\le1$, consists of tables admitting a stationary pair law $\Pi$ such that
\begin{equation}\label{eq:minor}
 \Pi P=\Pi,\qquad P^2(s,t)\ge\beta\Pi(t)\quad\text{for every pair }s,t.
\end{equation}
The certificate is existential and hidden. If it is nonunique, the stationary experiment includes its choice. In particular, the condition includes starting states outside $\supp\Pi$. Zero transition probabilities are allowed. In the stationary experiment initialization equals a qualifying hidden $\Pi$; in the arbitrary-start experiment $\mu$ is unrestricted.
\end{definition}

For an encoder family $\mathcal E$, define $\Risk_{\mathcal E,\mathrm{all}}$ as the infimum over encoder/decoder pairs independent of the source law, of the worst risk \eqref{eq:risk} over $\Dclass$ and all $\mu$; use $\Risk_{\mathcal E,\mathrm{stat}}$ for the stationary restriction. Public $K,M,Q,\beta$ may be known when choosing a pair. Source-independent common randomness is permitted, with every uniform architecture bound holding seedwise. The constructive predictors below are deterministic and do not read $\beta$ at a supplied budget.

\subsection{The complete cut and distinct resource coordinates}
The cut carries a vector $Z_d=Q^{-1}\floor{QF_d}$, finitely many exact boundary symbols $S$, and public metadata. Once the cut is made, the decoder may compute arbitrarily and output any real simplex vector. It may not read the history, clean numerator or denominator, an unquantized residual, another token's output, or a duplicate pre-cut state. This restriction is mathematical access control, not an uncharged physical isolation device.

Grid precision $h$ is not the complete arithmetic word length. A coordinate in $[0,1]$ on the closed grid has $Q+1$ labels and requires $h+1$ unsigned bits. A fixed bounded signed range requires $h+O(1)$ bits. Accumulators must additionally pay $\log M$, signs and full products. The scalar dimension $D$, number of physical heads, attention depth, total retained bits, and depth of an expanded Boolean/ReLU circuit are different quantities. A lower bound for one such coordinate is not a lower bound for all the others.

\subsection{Baseline bounded encoder models}
The \emph{original interface} retains the complete $9K+1$ scalar bank of four statistical heads, together with its declared exact query, endpoint and flag fields. Two additional paid singleton heads acquire local predecessor features. Its actual finite exponential evaluator, not ideal softmax, is part of the fixed encoder. The allowed decoder is optimized; the particular conservative decoder used for the upper bound is described in Section~\ref{sec:upper}.

For the \emph{one-summary local family}, fix integers $D,r\ge1$ and constants $0<a\le b<\infty$ independently of $M,\beta,Q$. It returns $D$ ratios
\begin{equation}\label{eq:localratio}
 F_d=\frac{\sum_{i=3}^{M+2} w_{di}(L_i,q(S))v_{di}(L_i,q(S))}
 {\sum_{i=3}^{M+2} w_{di}(L_i,q(S))},\qquad
 a\le w_{di}\le b,\quad |v_{di}|\le1.
\end{equation}
Here $L_i$ is a trailing window of fixed radius $r$, and $S$ contains the first and last $r$ symbols. The source-independent tables may vary with position and public parameters. Every coordinate pools all $M$ keys. Query-selective masks, nonlocal precomputed features and interior exact side statistics are excluded. $D$ counts scalar outputs; several coordinates may share a head and denominator. Arbitrary real tables are allowed for the information lower; the upper uses a specified finite construction.

The \emph{finite causal family} has a fixed architecture budget: depth, head count, dimensions, local radius, affine/ReLU operations per position, uniform state and coefficient bounds, score/value Lipschitz constants, and dense-mask fraction. History states depend only on causal local windows and public finite dyadic tables. All historical parameters are independent of future symbols. The exact boundary field, first and last three symbols in the basic version, may enter only the appended terminal initial query and the decoder. Public masks $\mathcal K_t\subseteq\{1,\ldots,t\}$ are independent of source symbols and boundary branches and satisfy $|\mathcal K_t|\ge ct$ for fixed $c>0$.

Queries, keys and values are affine maps, scores are bounded dot products with bounded public dyadic scale/bias, and positionwise updates/readout are fixed finite affine/ReLU directed acyclic graphs. Ideal states have a declared uniform range on all formal histories. Every public numerical datum is a supplied finite dyadic whose actual description length is charged. Layer normalization and unspecified real nonlinear evaluators are excluded from this finite subclass. Section~\ref{sec:arithmetic} specifies the integer evaluator used to implement each ideal encoder.

The density budget must permit $c\le1/3$. The architecture budget is assumed to accommodate one global block, four statistical heads, radius three, width $12K+4$, output dimension $D\ge9K+1$, affine row-sum bounds at least two, score bound at least two, state bound at least one and a sufficient score Lipschitz bound, for example $4K$. Larger fixed depths can add identity updates. Capacity does not certify arbitrary parameters: each selected network must obey the primitive bounds. Section~\ref{sec:binarycausal} states the smaller binary histogram capacity separately. Section~\ref{sec:depth} allows the number of stages to grow while fixing every per-stage primitive budget and making terminal initialization and arithmetic certification uniform. These are explicitly different comparison classes.

\section{Bounded-score risk and precision frontiers}\label{sec:main}
Write
\begin{equation}\label{eq:profile}
 \begin{aligned}
 \betaz&=\max\{\beta,M^{-1}\},&
 W(M,\beta,Q)&=\frac{1+\ln(1+Q/(M\betaz))}{Q},\\
 H(M,\beta,Q)&=\max\{M^{-1/2},W(M,\beta,Q)\}.&&
 \end{aligned}
\end{equation}
Constants below may depend on the fixed alphabet and architecture, but not on $M,\beta,Q$. They are not sharp joint-$K$ constants.

\begin{theorem}[Bounded finite information cut]\label{thm:main}
Assume $K\ge3$. For the original fixed interface,
\begin{equation}\label{eq:original}
 2^{-336}H\le\Risk_{\mathrm{orig},\mathrm{stat}}
 \le\Risk_{\mathrm{orig},\mathrm{all}}\le\min\{1,C_KH\},
 \qquad C_K=2^{16}K^{5/2}+2^{14}K^3+5320K^5.
\end{equation}
For the finite causal family of Section~\ref{sec:experiment}, evaluated at target precision $h+3$ by the integer compiler of Section~\ref{sec:arithmetic} and then rounded at resolution $Q$, there is $c_{\mathrm{arch}}>0$ such that
\begin{equation}\label{eq:causalmain}
 c_{\mathrm{arch}}H\le\Risk_{\mathrm{causal},\mathrm{stat}}
 \le\Risk_{\mathrm{causal},\mathrm{all}}\le\min\{1,C_KH\}.
\end{equation}
The same profile holds for \eqref{eq:localratio} when $D\ge9K+1$, $r\ge3$, $0<a\le1$ and $6\le b<\infty$, with a lower constant depending on $D,r,a,b$.
\end{theorem}

The original assertion is \unit{FR001 v1}, the local-family assertion is \unit{BF001 v3}, and the finite causal assertion is \unit{CF001 v1}. The underlying lower, finite arithmetic and upper steps are presented below; the precise dossier map is in Section~\ref{sec:map}. The three upper bounds do not identify the evaluated programs. In particular, the causal upper replaces the original head evaluator. Its deterministic accuracy envelope suffices; neither equality of old and new grid labels nor equality of selector decisions is claimed.

\begin{corollary}[Grid threshold]\label{cor:threshold}
Under Theorem~\ref{thm:main}, at fixed $K\ge3$ and architecture, the resolution required for statistical-order risk is, up to multiplicative constants,
\begin{equation}\label{eq:qstar}
 Q_* =\sqrt M\left[1+\ln\left(1+\frac1{\sqrt M\max\{\beta,M^{-1}\}}\right)\right].
\end{equation}
A dyadic budget within a constant factor of $Q_*$ attains a sufficiently large fixed multiple of $M^{-1/2}$. Attaining any specified fixed multiple requires $Q\gtrsim Q_*$, with the lower constant also depending on that multiplier. Thus
\[
 h_* =\tfrac12\log_2M+
 \log_2\left[1+\ln\left(1+\frac1{\sqrt M\max\{\beta,M^{-1}\}}\right)\right]+O(1).
\]
\end{corollary}
\begin{proof}
Let $s=\sqrt M$, $t=M\betaz$ and $f=1+\ln(1+s/t)$. For $Q=sf$, the inequality $1+sf/t\le f(1+s/t)$ gives $1+\ln(1+Q/t)\le f+\ln f\le2f$, so $W\le2/s$. Increasing $Q$ decreases $W$; dyadic rounding upward costs at most a factor two.
Conversely suppose $H\le A/s$ with fixed $A>0$, enlarged to $A\ge1$. Then $Q\ge s/A$ from $W\ge1/Q$. By concavity, $\ln(1+x/A)\ge A^{-1}\ln(1+x)$ for $A\ge1$. Therefore $1+\ln(1+Q/t)\ge f/A$, and $W\le A/s$ forces $Q\ge sf/A^2$. Apply the theorem's lower constant to a desired risk multiplier.
\end{proof}

For $\beta=0$ or $\beta=O(M^{-1})$, the threshold is $\sqrt M\log M$ and $h_*=\frac12\log_2M+\log_2\log M+O(1)$. If $\beta\sqrt M$ is bounded below, it is $\Theta(\sqrt M)$. The expression covers intermediate, length-dependent mixing promises without choosing a new decoder. Selecting a smaller budget from a known class promise is not learning an unknown $\beta$.

\section{Complete messages on stationary rare windows}\label{sec:lower}
The lower-bound proof has a source part and a representation part. The source is shared across local and causal encoders. The representation analysis must control \emph{every} retained coordinate and side field simultaneously.

\subsection{A legal stationary source with a common prior}
Use three symbols $W,C_0,C_1$. For $\beta\le1/16$, let
\[
 \alpha=\frac1{2M},\qquad \eta=\betaz-\alpha,\qquad
 p\in[1/4,3/4],\qquad p_0=1-p,\ p_1=p.
\]
Define a first-order kernel, embedded in the second-order model, by
\begin{align}\label{eq:witness}
 T(W,W)&=1-\alpha,& T(W,C_i)&=\alpha p_i,\notag\\
 T(C_i,W)&=\eta,& T(C_i,C_j)&=(1-\eta)p_j,\\
 \pi(W)&=\eta/\betaz,& \pi(C_i)&=\alpha p_i/\betaz.\notag
\end{align}
Unused-symbol rows equal $\pi$. Direct substitution gives $\pi T=\pi$ and $T(x,y)\ge\betaz\pi(y)$ for all symbols. Set $\theta_{uv}(x)=T(v,x)$ and $\Pi(u,v)=\pi(u)T(u,v)$. Then
\[
 P^2((u,v),(x,y))=T(v,x)T(x,y)\ge\betaz\Pi(x,y)\ge\beta\Pi(x,y),
\]
including off-support starting states. Initialize at this hidden stationary pair law.

For $Q\le M/4$, put
\begin{equation}\label{eq:jcut}
 J_* =\floor{\min\{(4\betaz)^{-1},M^2/(4Q^2)\}}.
\end{equation}
Here $4\le J_*\le M/4$. Let $E_j$ be the word event consisting of $M+2-j$ waiting symbols followed by $j$ active symbols, $2\le j\le J_*$. Its probability is independent of $p$ and obeys
\begin{equation}\label{eq:eventmass}
 e_j=\Pi(W,W)\alpha(1-\alpha)^{M-j}(1-\eta)^{j-1}\ge\frac1{32M}.
\end{equation}
Conditional on $E_j$, the active suffix is IID Bernoulli$(p)$ \emph{before} conditioning on endpoint information. The events are disjoint. One uniform prior on $p\in[1/4,3/4]$ is used across all $j$, and remains uniform conditional on each $E_j$. We may grant $j$ to the decoder without weakening a lower bound.

The actual next-symbol target on $E_j$ is
\begin{equation}\label{eq:target}
 \eta\delta_W+(1-\eta)((1-p)\delta_{C_0}+p\delta_{C_1}).
\end{equation}
The no-exit event in the observed suffix does not remove the next-step exit mass $\eta$. This distinction is essential for the stated loss.

\subsection{Local ratios and finite boundary branches}
For the local model \eqref{eq:localratio}, freeze a possible boundary value $s$ and its query while sampling all $j$ active bits. There are at most $H_b=3^{2r}$ such values on the witness. Write the numerator and denominator sums as $A_{ds}$ and $B_{ds}$ and define
\[
 m_{ds}(p)=\frac{\E_p A_{ds}}{\E_p B_{ds}}.
\]
This is a ratio of expectations, not an expectation of a ratio. Both denominators are at least $aM$. Replacing one coin changes at most $r$ summands, giving replacement bounds $2br$ and variances at most $2b^2r^2j$. Chebyshev and a union bound over all $2DH_b$ sums imply, with probability at least $9/10$, simultaneous deviations at most $T_b\sqrt j$, where $T_b=br\sqrt{40DH_b}$. Subtracting ratios then gives
\begin{equation}\label{eq:localtube}
 |F_{ds}-m_{ds}(p)|\le\frac{2T_b\sqrt j}{aM}.
\end{equation}
For a function of independent Bernoulli bits the exact polynomial identity
\begin{equation}\label{eq:derivative}
 \frac{d}{dp}\E_pf=\sum_i\E_p[f(B_i=1)-f(B_i=0)]
\end{equation}
implies $\|m_s'(p)\|_\infty\le4brj/(aM)$. Thus the whole bank lies near a short one-parameter curve, uniformly over arbitrary public position-dependent tables.

Choose a parameter mesh so that adjacent curve values differ by at most $1/Q$. Because $Q\sqrt j/M\le1/2$, every mesh point's tube intersects only a constant number of grid cells per coordinate. Counting all coordinates and the actual boundary label gives a fixed, $p$-independent set of at most
\begin{equation}\label{eq:labels}
 A_b(1+Qj/M),\qquad
 A_b=H_b(2+4br/a)\bigl(\ceil{8T_b/a+4}+1\bigr)^D
\end{equation}
complete good messages. The good event may depend on $p$; the containing alphabet may not. Freezing all possible boundary values first and selecting the actual one afterwards avoids a false assertion of conditional IID.

\begin{lemma}[Good-message Bayes bound]\label{lem:messages}
Let $p$ be uniform on an interval of length $\ell$. Suppose for every $p$ a good event has probability at least $9/10$, and every complete good message is in one fixed set of at most $R$ labels. Any possibly randomized scalar decoder has mean absolute error at least $\ell/(16R)$.
\end{lemma}
\begin{proof}
Fix the source-independent decoder seed. Radius-$\ell/(8R)$ neighborhoods of its at most $R$ good-message outputs cover at most one quarter of the prior interval. The joint probability of being good and outside those neighborhoods is at least $9/10-1/4$. On that event the error is at least $\ell/(8R)$; multiplying gives a bound stronger than the claim. Average over the seed. No posterior-uniformity assertion after the good event is used.
\end{proof}

By \eqref{eq:target}, $\TV(v,\theta)\ge(1-\eta)|v(C_1)/(1-\eta)-p|$. Applying the lemma, \eqref{eq:eventmass} and \eqref{eq:labels} yields
\begin{align}\label{eq:harmonic}
 \Risk&\ge\frac1{2048A_bM}\sum_{j=2}^{J_*}\frac1{1+Qj/M}\notag\\
 &\ge\frac1{8192A_bQ}\ln\left(1+\frac18\min\left\{\frac{Q}{M\betaz},\frac M Q\right\}\right).
\end{align}
The second cutoff $M/Q$ is real: it comes from the tube thickness in \eqref{eq:jcut}. It cannot simply be dropped.

\subsection{IID floor and completion of the resolution domain}\label{sec:completion}
On stationary binary IID sources, freezing the boundary gives an $O(1)$-length expectation curve and stochastic thickness $O(M^{-1/2})$. For $Q\le\sqrt M$, the same proof produces at most $A_IQ$ good complete messages and hence $\Risk\ge c_I/Q$. For $Q>\sqrt M$, and also outside the useful rare-window range, take Bernoulli parameters
\[
 p_\pm=\tfrac12\pm\frac1{16\sqrt{M+2}}.
\]
The history KL divergence is at most $1/8$; Pinsker and two-point testing give risk at least $1/(128\sqrt{M+2})\ge M^{-1/2}/256$, even with full history. These IID sources satisfy \eqref{eq:minor} for every $\beta\le1$. Therefore separate lower witnesses supply a statistical floor, a $1/Q$ floor and \eqref{eq:harmonic}; their maximum, not a sum of their losses, is available.

Here is the all-$Q$ completion used in both bounded-family theorems. If $\beta>1/16$, then $W\le Q^{-1}+16/M\le Q^{-1}+2M^{-1/2}$. If $Q>M/4$ or $Q>M^{3/4}$, $\ln(1+x)\le4x^{1/4}$ gives $W\le5M^{-1/2}$. In the remaining domain put $x=Q/(M\betaz)$ and $y=M/Q$. Since $x\le Q$ and $Q\le M^{3/4}$, $x>y$ implies $x\le y^3$. Thus
\[
 \ln(1+x)\le3\ln(1+\min\{x,y\})
 \le24\ln(1+\min\{x,y\}/8).
\]
The three lower bounds therefore control $H$ throughout the legal domain. This calculation neither extrapolates a thin-tube estimate nor imports an original-interface-only lower into a larger family.

\section{Causal depth and the correct influence envelope}\label{sec:causal}
Recursive causal features require a different representation argument from the local ratio calculation. Hold the appended initial query fixed within a boundary branch. Let the ideal state at row $t$ be $h_t$, and compare two formal histories with discrepancies $\delta_t=\|h_t-h'_t\|$. Assume a head has scores in $[-B,B]$, values bounded by $V$, score Lipschitz constant $S$, value Lipschitz constant $V_L$, and positionwise update Lipschitz constant $F_L$ for the sum norm. Write $H_a$ for the maximum heads per layer.

Along the line segment between two logit vectors, every weight is at most $C/t$, where $C=e^{2B}/c$ for masks of density $c$. The softmax differential is
\[
 d\alpha_i=\alpha_i\left(ds_i-\sum_k\alpha_kds_k\right),\qquad
 \sum_i|d\alpha_i|\le2\sum_i\alpha_i|ds_i|.
\]
Separating value and weight changes and integrating gives the conservative bound
\[
 \|o_t-o'_t\|\le2VCS\delta_t+
 \frac{C(V_L+2VS)}{t}\sum_{i\le t}\delta_i.
\]
The weight-change term is retained. Summing the heads and applying the update yields
\begin{equation}\label{eq:recurrence}
 \delta'_t\le a\delta_t+\frac b t\sum_{i\le t}\delta_i,
 \quad a=F_L(1+2H_aVCS),\quad b=F_LH_aC(V_L+2VS).
\end{equation}
No interpolation of hidden states or global Euclidean Lipschitzness of a dot product is needed; the primitive bounds apply on the declared reachable domain.

\begin{proposition}[Triangular influence]\label{prop:influence}
Let $A_{ti}=\ind\{i\le t\}/t$, $B_*=aI+bA$, and let the terminal row be $T=M+3$. For depth $L$, readout Lipschitz constant $G$ and initial symbol diameter $\Delta$, a replacement at $i<T$ changes the terminal vector by at most
\begin{equation}\label{eq:influence}
 d_i=G\Delta(B_*^L)_{Ti}
 \le\frac{G\Delta[(a+b)^L-a^L]}{T}
        [1+\ln(T/i)]^{L-1}.
\end{equation}
Moreover $\sum_{i<T}d_i\le C_0$ and $\sum_{i<T}d_i^2\le\Lambda/M$ for architecture constants, while $d_i\le\Gamma/M$ for $i\ge T/2$.
\end{proposition}
\begin{proof}
For $k\ge1$,
\[
 (A^k)_{ti}=\frac1t\sum_{i\le u_1\le\cdots\le u_{k-1}\le t}
 \frac1{u_1\cdots u_{k-1}}
 \le\frac1t\left(\sum_{u=i}^t\frac1u\right)^{k-1}
 \le\frac{[1+\ln(t/i)]^{k-1}}t.
\]
The empty product covers $k=1$. Binomial expansion of $B_*^L$ proves \eqref{eq:influence}. The row sum of $B_*^L$ is $(a+b)^L$, so the sum bound applies to its actual matrix entries, not merely their larger envelopes. For $m=2L-2$,
\[
 I_m=\int_0^1[1+\ln(1/x)]^m\,dx
     =\sum_{k=0}^m\binom m k k!<\infty.
\]
A decreasing-function Riemann sum gives $\Lambda=(G\Delta[(a+b)^L-a^L])^2I_m$. Restricting $i\ge T/2$ gives $\Gamma=G\Delta[(a+b)^L-a^L](1+\ln2)^{L-1}$. Radius-$r$ preprocessing multiplies the sum and late-influence constants by at most $r$ and the square-sum constant by at most $r^2$, since each source symbol affects at most $r$ initial states.
\end{proof}

A uniform $O(1/T)$ bound for every earlier token is false. With two zero-score scalar causal averaging layers, the map is $A^2$ and
\begin{equation}\label{eq:counterexample}
 (A^2)_{T1}=\frac1T\sum_{u=1}^T\frac1u=\frac{H_T}{T}.
\end{equation}
This is an admissible bounded one-head example. The theorem uses late-suffix influence and square summability, not the false uniform bound.

For a $j$-coin late suffix, Efron--Stein and \eqref{eq:derivative}, now applied to the entire nonlinear terminal output, give
\[
 \operatorname{Var}(F_d)\le\frac{\Gamma^2j}{2M^2},
 \qquad \| (\E_pF)'\|_\infty\le\frac{\Gamma j}{M}.
\]
For a fully IID history they give variance at most $\Lambda/(2M)$ and curve speed at most $C_0$. Thus the same tube count, common prior and domain completion apply, with new architecture constants. Unlike the local proof, the curve here is an expectation of the whole network output; no ratio-of-expectations substitution is made.

\subsection{Boundary information and temporal causality}
The proof can grant a stronger finite-side-information model: freeze each of finitely many boundary labels and allow branch-dependent tables and parameters. The uniform bounds must hold on \emph{all formal histories in each frozen branch}, including histories inconsistent with that branch's boundary value. Simultaneous branchwise concentration followed by selection of the actual label is valid, and its finite cardinality is counted.

That enlarged model need not be end-to-end causal. For example, assigning every initial history state the future bit $X_{M+2}$ through the branch label preserves causal attention edges but broadcasts future information to row one. The finite causal theorem excludes this behavior: historical tables, parameters and masks are independent of future boundary symbols, and the full boundary field enters only the appended terminal initial query. It is a subclass of the stronger lower-bound model. This scope restriction is not an extra statistical assumption or a claim that lower-triangular masks alone imply temporal causality.

\section{A constructive finite arithmetic bridge}\label{sec:arithmetic}
The influence argument concerns an ideal bounded architecture. Intermediate rounding is discontinuous and need not satisfy \eqref{eq:recurrence}. The bridge instead proves a uniform end-to-end approximation from explicit signed arithmetic.

\begin{theorem}[Finite evaluator]\label{thm:arithmetic}
For the finite-data affine/ReLU subclass in Section~\ref{sec:experiment}, an integer evaluator with fractional precision $P=4h+C_{\mathrm{arch}}$ satisfies
\[
 \|F_{\mathrm{fin}}(X)-F_{\mathrm{ideal}}(X)\|_\infty\le2^{-h}
\]
for every length and every formal history. Full online words additionally include $\log n$ accumulation guards and range/sign allowances. All supplied data, conversion, evaluation, movement and storage costs are charged. This assertion is \unit{FA001 v3}; it excludes arbitrary unevaluated Lipschitz maps and layer normalization.
\end{theorem}
\begin{proof}[Constructive proof]
Set $D_P=2^P$, $\varepsilon=D_P^{-1}$ and represent a signed real by integer $X$ with value $X/D_P$. Rounding is downward, including negative inputs. Convert every supplied dyadic parameter by exact signed division. For an affine coordinate, accumulate the complete signed integer
\[
 N_j=D_PB_j+\sum_k A_{jk}X_k,\qquad
 Y_j=\floor{N_j/D_P}.
\]
All products and the entire sum are retained before the single floor; ReLU is exact integer $\max(0,Y_j)$. With ideal input range $U$, discrepancy $\delta\le1$, and fan-in $d$, one affine step has error at most
\[
 \|A\|_\infty\delta+d\varepsilon(U+\delta)+2\varepsilon.
\]
Finite recursion yields computable constants for each fixed affine/ReLU block.

Interval bounds on the rounded-state cube of radius $R+1$ give an architecture-only integer $a\ge1$ bounding both ideal and finite logits in $[-a,a]$. Shift by $a$, which cancels in softmax. For the exact grid value $x=s_{\mathrm{fin}}+a\in[0,A]$, $A=2a$, let $S=D_Px$, $j_0=\ceil{8A^2}+2$ and $J=j_0+P$. Compute
\begin{equation}\label{eq:taylor}
 T_0=D_P,\qquad T_j=\floor{\frac{T_{j-1}S}{jD_P}},\qquad
 W=\sum_{j=0}^{J}T_j.
\end{equation}
The integer $H_a'=3^A$ bounds $e^x$. Each term-rounding error contributes at most $H_a'\varepsilon$ to the sum of all later terms. For $j\ge8A^2$, $A^j/j!\le2^{-j}$ by bounding the last half of the factorial. Hence
\begin{equation}\label{eq:exp}
 1\le W/D_P\le e^x,\qquad
 0\le e^x-W/D_P\le(H_a'J+1)\varepsilon.
\end{equation}
No exponential or logarithm is evaluated as an oracle.

For $m\ge1$ keys and signed value integers $V_i$, form
\begin{equation}\label{eq:signedmean}
 Z=\sum_iW_i\ge mD_P,\qquad N=\sum_iW_iV_i,\qquad
 O=\floor{N/Z},\qquad o_{\mathrm{fin}}=O/D_P.
\end{equation}
There is no extra division by $D_P$ inside $O$. The numerator uses full signed products with $\log m$ sum guards, so cancellation causes no invalid truncation. If finite and ideal values differ by at most $\eta_v$ and shifted weights by at most $\eta_w$, subtracting the two ratios gives
\[
 |o_{\mathrm{fin}}-o_{\mathrm{ideal}}|\le
 \varepsilon+\eta_v+2V\eta_w.
\]
Both weight sums are at least $m$ in real units, which removes dependence on key count. Combining this with the affine score bounds and \eqref{eq:exp} gives an error recursion
\[
 \delta_{\ell+1}\le a_*\delta_\ell+b_*(P+1)2^{-P}
\]
while $\delta_\ell\le1$. The readout and initialization obey the same finite recursion. Choose a computable $C_*\ge1$ dominating all resulting constants and select $C_{\mathrm{arch}}$ by finite search so that
\[
 C_*(C_{\mathrm{arch}}+5)2^{-C_{\mathrm{arch}}}\le1.
\]
For $P=4h+C_{\mathrm{arch}}$, this ensures $C_*(P+1)2^{-P}\le2^{-h}\le1$. Induction therefore keeps the actual states in the neighborhood used to derive the estimates, closing rather than assuming the bootstrap.
\end{proof}

If the finite output is within $E/Q$ of an ideal output, enlarging each good-message tube by $E/Q$ transfers the lower proof. With the compiler targeted at $h+3$, $E=1/8$; one may safely use $E=1$ in constants. An arbitrarily small error can cross a grid boundary, so this statement does not imply equality of rounded labels. It also does not say that the expanded finite arithmetic is a fixed-depth ordinary scalar Transformer.

\section{Terminal occupation and the constructive upper}\label{sec:upper}
This section identifies the inherited statistical inputs and the additional precision analysis. Let
\[
 J=\sum_{j=0}^{M-1}\ind\{Y_j=Y_M\},\qquad
 C_c=\sum_{j=0}^{M-1}\ind\{Y_j=Y_M,X_{j+3}=c\},\qquad
 H_c=\sum_{t=3}^{M+2}\ind\{X_t=c\}.
\]
These are proof variables, not decoder ports. Define the full-history add-one row $q_c=(C_c+1)/(J+K)$ and pooled row $b_c=(H_c+1)/(M+K)$.

\subsection{Inherited full-history statistical statements}
For precision about dependencies, put $n_0=M+1$, $\delta=1-\beta$ and
\begin{align*}
 i&=\min\{1,\sqrt{K/n_0}\},& a_0&=\min\{1,\sqrt{K^3/n_0}\},\\
 g&=\min\{1,\sqrt{K^5/n_0}\},& h_0&=\min\{1,\sqrt{K^4/n_0}\},\\
 F_0&=\max\{i,\min\{\delta,a_0\}\},&
 D_0&=\delta^{n_0}g,\quad T_0'=\delta^{M/2}h_0.
\end{align*}
The inherited \unit{R011} scales are
\begin{equation}\label{eq:inherited}
 \mathcal L=\max\{F_0,D_0,T_0'\},\qquad
 \mathcal H_{\mathrm{orig}}=\max\{a_0,D_0,T_0'\}\le g.
\end{equation}
Its \emph{ORIGINAL} assertion is $\sup\E\TV(q,\theta_{Y_M})\le2^{13}\mathcal H_{\mathrm{orig}}$. Its common, mixing-blind selector pools when
\[
 \TV(q,b)\le\rho(J),\qquad \rho(J)=\min\{g,64\sqrt{K/(J+K)}\},
\]
and has risk at most $2^{16}\mathcal L$. The underlying event-radius and rejection-restricted estimates, not just that scalar risk inequality, are used by the finite decoder. These statistical facts are inherited from the full \unit{R011}, \unit{R001}, \unit{S024} and \unit{S025} proof chain; they are not independently reproved in this article. In particular, no unconditional next-sample IID claim is substituted for them. Both scales are at most $K^{5/2}/\sqrt M$.

\subsection{From four bounded heads to a conservative decoder}
The compressed bank uses four score/value patterns at the terminal query $(a,b)$:
\begin{center}\small
\begin{tabular}{@{}lll@{}}\toprule
Head & Ideal score using supplied dyadics & Values\\\midrule
$P_1$ & $L_2\ind\{\mathrm{lag2}=a\}+L_2\ind\{\mathrm{lag1}=b\}$ & token, lag2\\
$P_2$ & $L_2\ind\{\mathrm{lag2}=a\}+L_3\ind\{\mathrm{lag1}=b\}$ & token, lag2\\
$T$ & $L_2\ind\{\mathrm{lag2}=b\}$ & lag1, initial flag vector\\
$0$ & $0$ & lag2, lag1, flag vector, flag\\\bottomrule
\end{tabular}
\end{center}
Here $L_2,L_3$ are finite downward approximations to $\ln2,\ln3$ and the flag vector at key three is $e_{X_2}$, zero elsewhere. The table has $2K+2K+2K+(3K+1)=9K+1$ outputs. All features are local bits; equalities are dot products of one-hots. Scores are at most $\ln6<2$, and values lie in $[0,1]$. The public mask is keys $3,\ldots,\min(t,M+2)$ at rows $t\ge3$, with self fallback at rows one and two. Its size is at least $t/3$ including the appended row. Only the terminal query uses future boundary information relative to historical positions.

The exact finite original weights lie in $[1,6]$. Their ratio formula verifies membership in the local family. For the causal-family upper, apply Theorem~\ref{thm:arithmetic} at target $h+3$ instead. The public log approximations are generated at
\[
 p_0=4h+4\ceil{\log_2(M+K+4)}+64.
\]
Their error perturbs each ideal head coordinate by at most $1/(8Q)$. Compiler error is at most $1/(8Q)$, and final downward rounding costs less than $1/Q$. Thus the new actual bank lies within $5/(4Q)<2/Q$ of the true-log identity bank. Positive Taylor weights and bit-valued projections also keep it on the closed $[0,1]$ grid.

The deterministic identity-and-error ledger in \unit{SRC} and \unit{NB002} then recovers grid frequencies $a_c,y_c\in[0,1]$ satisfying
\begin{equation}\label{eq:freq}
 |a_c-C_c/M|\le76/Q,\qquad |y_c-H_c/M|\le6/Q.
\end{equation}
All exact endpoint corrections and finite denominators are included in that ledger. This is the precise reusable interface; it does not require the old exponential recurrence once the head-error and range conditions are established.

Normalize positive weights $Ma_c+1$ and $My_c+1$ on the ordered $Q$-simplex grid, giving $\widehat q,\widehat b$. With $E_0=76KM/Q$, $A=J+K$ and $k=K/Q$, \eqref{eq:freq} implies
\[
 \TV(\widehat q,q)\le\min\{1,E_0/A\}+k,\qquad
 \TV(\widehat b,b)\le7k.
\]
Let $\widehat A=\sum_c(Ma_c+1)$, $A_-=\max\{K,\widehat A-E_0\}$, and
\[
 e_{\mathrm{obs}}=\min\{1,E_0/A_-\}+8k,\qquad
 \rho_-=\min\{g,64\sqrt{K/A_-}\}.
\]
The finite rule pools when $\TV(\widehat q,\widehat b)\le\min\{1,\rho_-+e_{\mathrm{obs}}\}$. The apparent square roots can be eliminated by the exact integer comparisons in \unit{NB002}; they are not runtime real oracles. Since $A_-\le A$, any history pooled by the exact selector is also pooled by the finite rule. Equality of decisions is unnecessary; the one-sided rejection containment is the relevant invariant.

Using the inherited event-radius and rejection estimates, \unit{NB002} establishes the pre-integration reduction
\begin{equation}\label{eq:preint}
 \operatorname{risk}\le\min\left\{1,2^{16}\mathcal L+
 3\E\min\left\{1,\frac{E_0}{J+K}\right\}
 +\Prob(J+K<4E_0)+\frac{9K}{Q}\right\}.
\end{equation}
This is a substantive inherited interface, not something inferred backwards from a final risk bound. The causal compiler satisfies its deterministic premises, so the same downstream map and reduction apply to the changed encoder.

\subsection{Dependence-safe occupation estimates}
For any homogeneous chain on $d$ states and arbitrary initialization, let $J_m$ count earlier visits to its actual state at time $m$. Then
\begin{equation}\label{eq:uc}
 \Prob(J_m<w)\le\min\{1,d^2w/(m+1)\},\qquad w\in\{1,2,\ldots\}.
\end{equation}
Indeed the largest such probability over initial states is nonincreasing in $m$ by terminal-suffix domination and the Markov property. On each deterministic path only the first $w$ occurrences of a state have fewer than $w$ earlier visits. Summing over time and then over the $d$ possible starting states bounds the sum of these maxima by $d^2w$. Monotonicity proves \eqref{eq:uc}. This is the inherited \unit{G/U} occupation argument, reproduced here because it is also used for the changed histogram encoder.

For $d=K^2$, layer cake at the exact integer threshold gives
\begin{equation}\label{eq:psi}
 \E\min\{1,A/(J_m+k)\}\le\Psi(K^4A/(m+1)),\quad k\ge1,
\end{equation}
where $\Psi(0)=0$, $\Psi(x)=x(1+\ln(1/x))$ on $(0,1)$ and $\Psi(x)=1$ on $[1,\infty)$. Specifically $J_m<A/t-k$ is either empty or $J_m<\ceil{A/t-k}$, and that ceiling is at most $A/t$. There is no independence assumption about the terminal query.

Under stationarity and $\beta>0$, \unit{IU001} proves the sharper bound
\begin{equation}\label{eq:stationaryocc}
 \E_\Pi\min\{1,A/(J_m+k)\}
 \le\frac{dA}{m}\left[17+\ln\left(1+\frac{480}{\beta A}\right)\right],
 \quad d=|\supp\Pi|\le K^2.
\end{equation}
The proof reverses the chain conditional on its actual endpoint $s$. The reversed kernel on positive stationary support satisfies the same two-step minorization. If $p=\Pi(s)$ and $\tau_w$ is its $w$th positive return, the stopped occupation-vector argument gives $\E_s\tau_w=w/p$ and $\Prob_s(J_m<w)\le w/(mp)$. A tilted singular-operator comparison, valid without reversibility, gives
\[
 \Prob_s(J_m<mp/16)\le5e^{-\beta mp/96}.
\]
Integrating the minimum of this tail and $A/(mpt)$, then averaging over $s$, yields \eqref{eq:stationaryocc}. The factor $p$ cancels only after the conditional estimate has been proved. The full operator and return proof is retained in \unit{IU001}, rather than hidden behind an informal mixing-concentration citation.

\subsection{Exact residual suffix laws}
A TV burn error alone would leave a nonvanishing additive term as $Q\to\infty$. The all-$Q$ upper instead uses an exact path mixture. Let $u=\floor{M/2}$, $m=M-u$, $r=\floor{M/4}$ and, for $0<\beta<1$,
\[
 R=\frac{P^2-\beta\mathbf\Pi}{1-\beta},\qquad
 \mathbf\Pi(s,t)=\Pi(t),\qquad \epsilon=(1-\beta)^r.
\]
All-state minorization makes $R$ stochastic even off support. Since $R\mathbf\Pi=\mathbf\Pi R=\mathbf\Pi$, for every initial $\mu$,
\begin{equation}\label{eq:mixture}
 \mu P^u=(1-\epsilon)\Pi+\epsilon\nu,
 \qquad \nu=\mu R^rP^{u-2r}.
\end{equation}
Apply the same original-$P$ path-extension kernel to both components: the \emph{entire} suffix law is a mixture of a stationary $P$ trajectory and an ordinary $P$ trajectory from $\nu$. The residual suffix does not evolve under $R$. At $\beta=1$ the stationary component is exact and the residual is omitted.

Full-history visits dominate suffix visits pathwise. Apply \eqref{eq:stationaryocc} to the stationary component and \eqref{eq:psi} to the residual component in \eqref{eq:preint}. The result from \unit{AU001} is
\[
 \operatorname{risk}\le\min\left\{1,2^{16}\mathcal L+
 \frac{2^{14}K^3}{Q}\left[1+\ln\left(1+\frac Q{\beta KM}\right)\right]
 +7\epsilon\Psi(152K^5/Q)\right\}.
\]
For $\beta M=t\ge1$, $\epsilon\le e e^{-t/4}$ and
\[
 \epsilon(1+\ln Q)\le1+\ln(1+Q/t)+4.
\]
This bounds the residual by $5320K^5W$. If $\beta<1/M$, the universal occupation bound directly yields the same envelope. Hence
\[
 \operatorname{risk}\le\min\{1,2^{16}\mathcal L+(2^{14}K^3+5320K^5)W\}
 \le\min\{1,C_KH\}.
\]
The suffix, reverse chain and residual kernel are analysis objects. The predictor does not compute them or receive another history port. Together with the lower proof and stationary inclusion this completes the theorem's stated modular argument.

\section{Weighted inventories and the effective score budget}\label{sec:score}
Bounded-range families have a genuine representation restriction. We now change the finite evaluator and retain more coordinates. The construction and upper bounds are \unit{SB001 v2}; the sharp weight-ratio lower is \unit{SH001 v1}. Neither reinterprets the old exponential program.

Two paid singleton-copy heads acquire the previous two one-hots. Paid local Boolean gates form pair keys $p_t(a,b)$ and query-independent one-hot triple values $v_t(a,b,c)$. A dense statistical head has key/query dimension $K^2$ and retains all $K^3$ coordinates. Its exact pair dot product selects the integer weight $\kappa$ on matches and $1$ otherwise, where $\kappa\ge1$ is public. Its effective natural logits are $0,\ln\kappa$; no logarithm or exponential is executed. All $M$ terminal keys have positive weight.

At every row compute exact weighted numerators and denominator, then retain only
\[
 z_{abc}=\floor{Q A_{abc}/D}.
\]
The decoder selects the $K$ coordinates for the exact terminal pair and normalizes their sum. If all selected floors vanish, it uses the uniform vector. Ordered simplex rounding gives its final $Q$-grid output. The clean $A_{abc},D$, counts and unquantized block mass are not decoder inputs.

\subsection{Deterministic error including the zero-floor branch}
Let $J$ and $C_c$ denote the earlier occurrences and successors of the terminal pair. For $J>0$, the exact selected block has mass
\begin{equation}\label{eq:blockmass}
 r=\frac{\kappa J}{M+(\kappa-1)J}
\end{equation}
and conditional composition $C/J$. Nonmatches occupy other coordinates rather than contaminating this composition. If nonnegative $y_c$ sum to $r$ and $f_c=Q^{-1}\floor{Qy_c}$, their lost mass is less than $K/Q$. When $\sum f_c>0$, subtracting normalized vectors gives $\TV(f/\sum f,y/r)\le K/(Qr)$. If all floors vanish, $r<K/Q$, so the same capped bound covers any simplex fallback.

Comparison to add-one smoothing and final ordered rounding therefore gives, on every history including $J=0$,
\begin{equation}\label{eq:inventoryerr}
 \TV(T_\kappa,q)\le\frac K{J+K}+
 \min\left\{1,\frac{K(K+1)M}{\kappa Q(J+K)}\right\}
 +\frac{2K-1}{Q}.
\end{equation}
The inequality uses $1/r=(\kappa-1)/\kappa+M/(\kappa J)$ and $J+K\le(K+1)J$ when $J\ge1$. The first term equals one at $J=0$. No positive mass surviving the grid is assumed.

\begin{theorem}[Amplified triple bank]\label{thm:scoreupper}
Let $D_K=K^5(K+1)$. For arbitrary initialization and every mixing promise,
\begin{equation}\label{eq:scoreupper}
 \operatorname{risk}(T_\kappa)\le\min\left\{1,
 2^{13}\mathcal H_{\mathrm{orig}}+\Psi\left(\frac{K^5}{M+1}\right)
 +\Psi\left(\frac{D_K}{\kappa Q}\right)+\frac{2K-1}{Q}\right\}.
\end{equation}
With the integer, mixing-blind choices
\[
 L_M=\bl(M+1),\quad \kappa=L_M,\quad
 h=\max\{4,\ceil{L_M/2}\},\quad Q=2^h,
\]
the risk is at most
\begin{equation}\label{eq:integerpoint}
 \min\left\{1,\frac{(2^{13}+2)K^{5/2}+3D_K+2K-1}{\sqrt{M+1}}\right\}.
\end{equation}
Thus the effective logit range is $\ln\ln(M+1)+O(1)$ at $Q\asymp\sqrt M$.
\end{theorem}
\begin{proof}
Apply \eqref{eq:psi} to the first two terms in \eqref{eq:inventoryerr}, and use the inherited \unit{R011 ORIGINAL} risk bound. This proves \eqref{eq:scoreupper}. At the integer schedule, $\sqrt{M+1}\le Q\le2\sqrt{M+1}$ and $\ln(1+Q)\le L_M$. For $D\ge1,z>0$,
\[
 \Psi(D/z)\le\frac{D[1+\ln(1+z)]}{z},\qquad \Psi(x)\le2\sqrt x.
\]
Taking $z=L_MQ$ gives $\Psi(D_K/(L_MQ))\le3D_K/Q$, which proves \eqref{eq:integerpoint}. Budget generation needs only bitlength, integer ceiling and a power of two.
\end{proof}

\subsection{A rescaled tube with uniform weight-ratio constants}
Fix $K\ge3$ and a one-block local family with at most $D$ coordinates, trailing radius $r$, actual weights in $[1,\kappa]$, values in $[-1,1]$, and only the boundary side information of \eqref{eq:localratio}. Here $\kappa\ge1$ may be real; $D,r$ are fixed independently of $M,\beta,Q,\kappa$. The lower permits arbitrary real public positional tables and source-independent shared randomness. It does not include recursively global features or query-selected masks. Smaller banks can be conceptually zero-padded for this information argument without changing a physical implementation.

The earlier \unit{SB001} substitution in \unit{BA001} gives the valid lower $\Omega(\ln M/(\kappa^{D+1}\sqrt M))$ in weak mixing at $Q\asymp\sqrt M$. It loses a power of $\kappa$ through the number of grid cells across a thickened tube. The following refinement changes the suffix cutoff instead of treating that old constant as uniform.

\begin{proposition}[Rescaled rare-window lower]\label{prop:rescaled}
Put $z=\kappa Q$, $H_b=3^{2r}$, $T_b=r\sqrt{40DH_b}$ and
\[
 N_b=\ceil{8T_b+4}+1,\qquad
 A_b^*=H_b(2+4r)N_b^D,\qquad c_b=(8192A_b^*)^{-1}.
\]
For $0\le\beta\le1/16$ and $z\le M/4$, the stationary minimax risk, optimized over this entire family, satisfies
\begin{equation}\label{eq:rescaled}
 \Risk_{\kappa,\mathrm{stat}}\ge
 \frac{c_b}{z}\ln\left(1+\frac18\min\left\{\frac{z}{M\betaz},\frac Mz\right\}\right).
\end{equation}
The constant depends only on $D,r$, not on $\kappa,Q,M,\beta,K$ or the tables. This is \unit{SH001 v1}; it is not a matching all-$Q$ two-resource risk formula.
\end{proposition}
\begin{proof}
Use the stationary witness \eqref{eq:witness}, the same uniform prior and the events $E_j$, but replace \eqref{eq:jcut} by
\[
 J_*=\floor{\min\{(4\betaz)^{-1},M^2/(4z^2)\}}.
\]
Both terms are at least four, so $J_*\ge4$, $J_*$ is at least half the unfloored minimum, and $z\sqrt{J_*}/M\le1/2$. Freeze every boundary query before sampling the suffix. A coin changes at most $r$ numerator or denominator terms, by at most $2\kappa r$ in total. The simultaneous sum deviations are therefore at most $\kappa T_b\sqrt j$ with probability at least $9/10$. Denominators are at least $M$, giving
\[
 |F_{ds}-m_{ds}(p)|\le2\kappa T_b\sqrt j/M,\qquad
 \|m_s'(p)\|_\infty\le4\kappa rj/M.
\]
After scaling by $Q$, the tube thickness is at most $T_b$, independent of $\kappa$. A parameter mesh with gaps at most $M/(4rzj)$ has at most $(2+4r)(1+zj/M)$ points and one grid unit of approximation error. Rounding and the boundary label give at most $A_b^*(1+zj/M)$ complete good messages. Lemma~\ref{lem:messages} and \eqref{eq:eventmass} now yield
\[
 \Risk_{\kappa,\mathrm{stat}}\ge
 \frac1{2048A_b^*M}\sum_{j=2}^{J_*}(1+zj/M)^{-1}.
\]
The harmonic comparison from Section~\ref{sec:lower}, with $z$ in place of $Q$, proves \eqref{eq:rescaled}. The prior is unchanged across events and independent seeds. The constant is uniform because $\kappa$ shortens the admissible window and enters the mesh length through $z$; it no longer enlarges the grid-box cross-section.
\end{proof}

\subsection{The least feasible budget in a square-root grid band}
Fix $0<c_-\le c_+<\infty$, a legal dyadic $Q$ with $c_-\sqrt M\le Q\le c_+\sqrt M$, and a target multiplier $C>0$. For $a\in\{\mathrm{stat},\mathrm{all}\}$ define
\begin{equation}\label{eq:kappamin}
 \kappa_{\min,a}(M,\beta,Q;C)
 =\inf\{\kappa\in[1,\infty):\Risk_{\kappa,a}\le C/\sqrt M\},
\end{equation}
with $\inf\varnothing=+\infty$. The envelope budget $\kappa$ is real and the lower allows real tables; a feasible upper below uses actual integer weights. This definition assumes neither attainment of the infimum nor a noninteger implementation.

The unchanged \unit{SB001} upper is more precise than \eqref{eq:scoreupper}. Put
\[
 t=M\betaz,\quad G(z,t)=1+\ln(1+z/t),\quad
 C'_K=48K^3(K+1)+10D_K.
\]
Then its same paid triple bank obeys
\begin{equation}\label{eq:scoremix}
 \operatorname{risk}(T_\kappa)\le\min\left\{1,
 2^{13}\mathcal H_{\mathrm{orig}}+\Psi\left(\frac{K^5}{M+1}\right)
 +C'_K\frac{G(\kappa Q,t)}{\kappa Q}+\frac{2K-1}{Q}\right\}.
\end{equation}

\begin{theorem}[Sharp weight-ratio order in the fixed band]\label{thm:kappa}
Let $K\ge3$, $D\ge K^3$ and $r\ge3$ be fixed. Set
\[
 f=1+\ln(1+Q/(M\betaz)),\qquad
 C_{\mathrm{suff}}=(2^{13}+2)K^{5/2}+(2C'_K+2K-1)/c_-.
\]
For every fixed $C\ge C_{\mathrm{suff}}$, uniformly in $M\ge64$, $\beta\in[0,1]$ and all legal $Q$ in the stated band,
\begin{equation}\label{eq:kappasharp}
 c_*f\le\kappa_{\min,\mathrm{stat}}
 \le\kappa_{\min,\mathrm{all}}\le\ceil f\le2f,
\end{equation}
where $c_*>0$ depends on $D,r,c_+,C$. The necessary bound holds for every fixed $C>0$; feasibility at arbitrarily tiny multipliers is not claimed.
\end{theorem}
\begin{proof}
For sufficiency choose the public integer $\kappa=\ceil f$. Since $1+\kappa Q/t\le\kappa(1+Q/t)$,
\[
 G(\kappa Q,t)/\kappa\le(f+\ln\kappa)/\kappa\le2.
\]
Use $\mathcal H_{\mathrm{orig}}\le K^{5/2}/\sqrt{M+1}$, $\Psi(x)\le2\sqrt x$ and $Q\ge c_-\sqrt M$ in \eqref{eq:scoremix}. This proves arbitrary-start feasibility with $C_{\mathrm{suff}}$. Its all-zero-block and $J=0$ branches remain covered by \eqref{eq:inventoryerr}.

For necessity write $x=Q/(M\betaz)$, so $x\le Q$ and $Q^2/M\le c_+^2$. Choose a finite $F\ge2(1+\ln8)$, depending only on $c_+$, such that for every $s\ge F$,
\[
 \frac{e^{s-1}-1}{8c_+^2s}\ge e^{s/2}-1.
\]
Such a choice exists by exponential domination. Exhaust the cases. If $\beta>1/16$, then $x\le2c_+$ and $\kappa\ge1\ge f/(1+2c_+)$. If $z>M/4$, then $\kappa>\sqrt M/(4c_+)$, whereas $f\le(1/8+c_+)\sqrt M$; this gives a fixed multiple of $f$ without extrapolating the tube. In the remaining domain \eqref{eq:rescaled} applies. If $f\le F$, use $\kappa\ge f/F$; if $\kappa\ge f$, there is nothing to prove.

Otherwise $f>F$ and $1\le\kappa<f$. The two logarithmic branches satisfy
\[
 \ln(1+\kappa x/8)\ge f-1-\ln8\ge f/2,\qquad
 \ln(1+M/(8\kappa Q))\ge f/2.
\]
For the second, use $M/(\kappa Q)\ge x/(\kappa c_+^2)>x/(fc_+^2)$ and $x=e^{f-1}-1$. Thus \eqref{eq:rescaled} and risk at most $C/\sqrt M$ force $\kappa\ge c_bf/(2Cc_+)$. Taking the minimum of the five positive constants
\[
 1,\quad (1+2c_+)^{-1},\quad
 [4c_+(1/8+c_+)]^{-1},\quad F^{-1},\quad c_b/(2Cc_+)
\]
proves uniform necessity. It applies to every feasible real budget, hence its infimum. Stationary inclusion gives the middle inequality in \eqref{eq:kappasharp}.
\end{proof}

When $\beta\le1/M$, $f\asymp\ln M$ on this band, so $\kappa_{\min,a}=\Theta(\ln M)$ and $\ln\kappa_{\min,a}=\ln\ln M+O(1)$. For bounded $f$, the meaningful uniform range comparison is $1+\ln\kappa_{\min,a}=\Theta(1+\ln f)$, not a positive lower on $\ln\kappa_{\min,a}$. The integer upper changes no program or cost ledger. Specifying $\ceil f$ from a known class promise is not learning $\beta$ or evaluating a free transcendental primitive: an implemented selection rule must pay its supplied-parameter and arithmetic costs. The mixing-blind bitlength schedule in Theorem~\ref{thm:scoreupper} remains available. Bands containing no legal dyadic $Q$ have no instance at that $M$.

\subsection{Uniform attention and representation costs}
At $\kappa=1$, the construction is ordinary uniform attention over the full triple histogram, followed by conditional block normalization. It is a strong elementary baseline, not a new counting mechanism.
\begin{corollary}[Uniform full-histogram synthesis]\label{cor:hist}
At fixed $K\ge3$, its best post-grid decoder has stationary and arbitrary-start risk of order $H(M,\beta,Q)$, uniformly over the legal $\beta,Q$.
\end{corollary}
\begin{proof}
For the upper, set $\kappa=1$ in \eqref{eq:scoremix}: the first two terms are $O_K(M^{-1/2})$, since $\Psi(x)\le2\sqrt x$, and the remaining terms are $O_K(W)$. For the lower, this exact $K^3$ bank belongs to \eqref{eq:localratio} with $a=b=1,r=3$; its decoder risk obeys the local-family $1/Q$, statistical and rare-window lower bounds, combined in Section~\ref{sec:completion} into $\Omega_K(H)$. This is a synthesis of \unit{SB001} and \unit{BA001}, not a separately audited source unit.
\end{proof}

The full bank pays $K^3$ coordinates instead of the compressed bank's $9K+1$, so fixed-$K$ rate equality does not imply equal dimension or costs. The earlier \unit{GW001 v3} successor-only construction with weights $1,M,M^2$ still gives a valid sufficient point with effective range $2\ln M$. The old \unit{SB001} lower also remains true; the rescaled proof sharpens it rather than refuting it. The fixed-band theorem above has no exact threshold constant and no implication for arbitrary deep or growing architectures. Section~\ref{sec:joint} extends the one-block analysis to all legal joint weight/grid regimes using an additional uniform IID argument.

\section{The joint weight-ratio and grid frontier}\label{sec:joint}
The fixed-band threshold of Section~\ref{sec:score} does not by itself determine the risk when both resources vary. In particular, the earlier IID tube constant depended on the weight envelope. The next argument proves that even an arbitrarily large finite weight ratio cannot remove the final-grid penalty. Together with the rescaled rare-window lower and the unchanged triple-bank upper, this closes the joint frontier in the one-summary-block family. It does not enlarge the bounded causal family of Theorem~\ref{thm:main}.

\subsection{A grid floor uniform in the weight envelope}
\begin{proposition}[Uniform IID grid floor]\label{prop:iiduniform}
Fix $D,r\ge1$ in the local-ratio model \eqref{eq:localratio}, with actual weights in $[1,\kappa]$. For all finite real $\kappa,Q\ge1$ and $M\ge64$,
\[
 \Risk_{\kappa,\mathrm{stat}}\ge c_I(D,r)/Q
\]
even on stationary binary IID sources. The constant is independent of the weight envelope, alphabet size and mixing promise. The decoder may still return arbitrary real probabilities. This is \unit{UI001 v1}.
\end{proposition}
\begin{proof}
Use IID Bernoulli$(p)$ observations on two designated symbols, and define every second-order row, including unsupported rows, to emit that same law. Then $P^2(s,t)=\Pi(t)$ on every state, so every $\beta\in[0,1]$ is allowed and the actual terminal target is always $p$. Fix and grant an encoder seed. The following prior and constants do not depend on that seed.

Write
\[
 H_0=3^{2r},\quad F_0=DH_0,\quad q_0=4^{-r},\quad L_0=rF_0+2r,
\]
\[
 J_0=(2rL_0+1)/q_0,\qquad V_0=4r^2J_0/q_0,
\]
and choose an integer
\[
 C_0\ge\max\{1,\ 400F_0(2r-1)/q_0^2,\ 3200F_0V_0\}.
\]
Put $m=\ceil{C_0Q^2}$. If $M<m$, the full-history two-point bound of Section~\ref{sec:completion} gives risk at least
$1/(128\sqrt{M+2})\ge1/(128\sqrt{C_0+3}\,Q)$.
This branch already handles a locality as large as the history.

Suppose $M\ge m$ and give $p$ the uniform prior on $I=[1/4,3/4]$. Freeze every possible boundary branch $s$ before sampling. For each coordinate, branch and position, let $a_{dsi}$ be the largest weight over the attainable binary patterns in its trailing window. Select the $m$ largest such maxima for each $(d,s)$, with public tie-breaking. Reveal all bits in those windows and all boundary bits. The revealed position set $U$ is fixed independently of the sample and $p$; for $N=|U|$,
\[
 m\le N\le rF_0m+2r\le L_0m.
\]
This is a proof-only oracle, not a new encoder feature. Write $B=X_U$.

A chosen maximizing pattern has probability at least $q_0$. Each selected window overlaps at most $2r-2$ others, so the variance of the count of realized maximizing patterns is at most $(2r-1)m$. Chebyshev and a union over the $F_0$ branches give a $B$-measurable event $E$, with probability at least $.99$ for every $p$, on which at least $q_0m/2$ selected patterns realize their maxima in every branch. If $a_*$ is the largest unselected maximum in one branch, its now-known selected weight mass is at least $(q_0m/2)a_*$.

Condition on $B\in E$ and fixed $p$. For a numerator $A$ and denominator $D_n$, let $\mu_A=\E[A\mid B]$, $\mu=\E[D_n\mid B]$ and $h=\mu_A/\mu$. At most $rN$ unselected windows intersect $U$, and their total maxima are bounded by $rL_0ma_*$. Do not assign probability $q_0$ to their old maximizing patterns: those patterns may have become impossible after the reveal. Each disjoint window, however, retains its IID law and has expected weight at least $q_0a_i$. Charging the intersecting windows to the selected mass gives
\[
 \sum_{\mathrm{unselected}}a_i\le J_0\mu,\qquad
 a_*\le\frac{2\mu}{q_0m}.
\]
Efron--Stein on the remaining independent bits, with at most $r$ affected summands per bit, therefore yields
\[
 \Var(A\mid B),\ \Var(D_n\mid B)
 \le2r^2\sum a_i^2\le V_0\mu^2/m.
\]
A simultaneous Chebyshev event of conditional probability at least $.99$ bounds both deviations by $\mu/(4Q)$. Exact ratio subtraction then gives
\[
 |A/D_n-h|\le 2/(3Q)<1/Q.
\]
For every remaining local pattern the absolute logarithmic derivative of its Bernoulli probability is at most $4r$. Positivity of weights gives
$|\mu_A'|,|\mu'|\le4r\mu$, hence $|h'|\le8r$, independently of $\kappa$.

The oracle posterior is not uniform. If $B$ contains $t_B$ ones and $y=t_B/N$, its density on $I$ is proportional to $p^{t_B}(1-p)^{N-t_B}$. Put
\[
 I_B=I\cap[y-3/\sqrt N,y+3/\sqrt N].
\]
Hoeffding followed by Markov shows that, except for a set of $B$ of probability at most $200e^{-18}$, this interval has posterior mass at least $.99$. Call this event $T$. On $T$ its posterior density is at most $128\sqrt N$: for the clamped mode $p_0$, the log density $\ell$ satisfies
$|\ell'(p_0)|\le16\sqrt N$ and $|\ell''|\le16N$. Integrating for distance $1/(64\sqrt N)$ in a direction remaining inside $I$ gives a normalizer at least $e^{\ell(p_0)}/(128\sqrt N)$. This also handles a mode on a truncation boundary. Consequently the density is at most
\[
 P_0Q,\qquad P_0=128\sqrt{L_0(C_0+1)}.
\]

For $B\in E\cap T$, an economical mesh of $I_B$ with spacing at most $1/(8rQ)$ has at most $2+48r$ points, because $N\ge Q^2$. The derivative bound, random tube and downward grid together put every good complete bank in at most
\[
 R_0=(2+48r)9^D
\]
labels. There is no missing boundary factor: $B$ fixes the actual side field. The joint posterior probability of lying in $I_B$ and being in the tube is at least $.98$. For a fixed decoder seed, neighborhoods of radius
$\varepsilon=(8R_0P_0Q)^{-1}$ about all $R_0$ possible scalar outputs have posterior mass at most $1/4$. The conditional expected error is thus at least $1/(16R_0P_0Q)$. Since $\Prob(E\cap T)>.98$, averaging over $B$ and both seeds proves the result with
\[
 c_I=\min\{(128\sqrt{C_0+3})^{-1},(32R_0P_0)^{-1}\}.
\]
The fixed reveal and its posterior are both essential. Treating the reveal as free model information, or retaining a uniform prior after it, would not prove this statement.
\end{proof}

\subsection{All joint regimes and both initializations}
For $t\ge1$ define
\[
 \mathcal W(z,t)=\frac{1+\ln(1+z/t)}{z},\qquad
 \mathcal F(M,\beta,\kappa,Q)
 =\max\left\{M^{-1/2},\,Q^{-1},\,\mathcal W(\kappa Q,M\betaz)\right\}.
\]
Here $\kappa$ is a weight envelope, not an additional exact message. The following theorem permits arbitrary real tables inside that envelope; its upper witness uses only integer weights.
\begin{theorem}[Complete one-block frontier]\label{thm:joint}
Fix $K\ge3$, $D\ge K^3$ and $r\ge3$. In the one-block local positive-ratio family, for every integer $\kappa\ge1$, $M\ge64$, $\beta\in[0,1]$ and legal dyadic $Q\ge16$,
\begin{equation}\label{eq:joint}
 c\,\mathcal F\le\Risk_{\kappa,\mathrm{stat}}
 \le\Risk_{\kappa,\mathrm{all}}
 \le\operatorname{risk}_{\mathrm{all}}(T_\kappa)
 \le\min\{1,C_K^{\mathrm{joint}}\mathcal F\}.
\end{equation}
One may take
\[
 c=\min\{(256\cdot5)^{-1},c_I/2,c_b/48\},\qquad
 C_K^{\mathrm{joint}}=(2^{13}+2)K^{5/2}+C'_K+2K-1.
\]
The lower constants depend only on the fixed $D,r$, not on $\kappa$. This is \unit{TR001 v1}; it is not a theorem for arbitrary deep or growing architectures.
\end{theorem}
\begin{proof}
The statistical floor from Section~\ref{sec:completion} gives $\Risk\ge M^{-1/2}/256$, and Proposition~\ref{prop:iiduniform} gives $c_I/Q$. Proposition~\ref{prop:rescaled} supplies $c_bL$ only for $\beta\le1/16$ and $z=\kappa Q\le M/4$, where
\[
 L=z^{-1}\ln(1+\min\{z/t,M/z\}/8),\qquad t=M\betaz.
\]
These bounds may use different witnesses; take their maximum, not a sum of losses.

The completion is elementary but must be uniform in $z$. The function $\mathcal W$ decreases in $z$ and $t$. If $\beta>1/16$,
$\mathcal W\le Q^{-1}+16/M\le Q^{-1}+2M^{-1/2}$.
If $z>M/4$, then
\[
 \mathcal W(z,t)\le4[1+\ln(1+M/4)]/M<5M^{-1/2}.
\]
If $z>M^{3/4}$, then
\[
 \mathcal W(z,t)\sqrt M
 \le\frac{1+\ln(1+z)}{z^{1/3}}
 \le1+\ln2+3/e<3.
\]
On the remaining domain, set $x=z/t$, $y=M/z$. Then $x\le z\le y^3$ and $y\ge4$. Splitting at $x=y$ and using log concavity gives
\[
 \ln(1+x)\le24\ln(1+\min\{x,y\}/8),\qquad
 \mathcal W\le Q^{-1}+24L.
\]
The separate floors and the truncated lower therefore imply the displayed $c\mathcal F$ everywhere, including equality at the branch boundaries.

For the upper, apply \eqref{eq:scoremix} and
$\mathcal H_{\mathrm{orig}}\le K^{5/2}/\sqrt{M+1}$,
$\Psi(u)\le2\sqrt u$. The complete triple bank is in this family because $D\ge K^3$, $r\ge3$, all keys have positive weights, and its query uses only the permitted terminal pair. The zero-block fallback, $J=0$ case and post-grid normalization remain unchanged. Stationary inclusion gives the middle inequality. All full-history statistical input and whole-suffix mixture dependencies remain exactly those identified in Section~\ref{sec:upper}.
\end{proof}

\begin{corollary}[The two simultaneous statistical-order budgets]\label{cor:jointthreshold}
At fixed $K,D,r$, statistical-order risk requires and, up to constant factors, is achieved by
\begin{equation}\label{eq:jointthreshold}
 Q\gtrsim\sqrt M,\qquad
 \kappa Q\gtrsim Z_*,
 \quad
 Z_*=\sqrt M\left[1+\ln\left(1+\frac1{\sqrt M\betaz}\right)\right].
\end{equation}
Necessity holds for every prescribed fixed target multiplier; sufficiency asserts an explicit sufficiently large multiplier. The quantity $Z_*$ equals the earlier numerical threshold $Q_*$, but is now a \emph{product} threshold, not a replacement for the separate grid requirement.
\end{corollary}
\begin{proof}
Put $a=\sqrt M$, $u=a/(M\betaz)$, $f_*=1+\ln(1+u)$ and $v=\kappa Q/a$. If risk is at most $C_0/a$, let $B=\max\{1,C_0/c\}$. Theorem~\ref{thm:joint} gives $Q\ge a/B$ and
$[1+\ln(1+vu)]/v\le B$.
For $v\ge1$ the numerator is at least $f_*$, so $v\ge f_*/B$.
For $v<1$, concavity gives $1+\ln(1+vu)\ge vf_*$, hence $f_*\le B$; the numerator is also at least one, so $v\ge1/B\ge f_*/B^2$. Thus $\kappa Q\ge Z_*/B^2$ in every case.

Conversely, $1+f_*u\le f_*(1+u)$ implies
$\mathcal W(af_*,M\betaz)\le2/a$.
If $Q\ge a$ and $\kappa Q\ge Z_*$, the unchanged program has risk at most $2C_K^{\mathrm{joint}}/a$. Fixed multiplicative relaxations of either budget change only this target constant. Legal choices always exist: take
$Q=2^{\max\{4,\ceil{\log_2\sqrt M}\}}$ and
$\kappa=\max\{1,\ceil{Z_*/Q}\}$.
\end{proof}

At a fixed square-root grid band, \eqref{eq:jointthreshold} recovers the sharp weight-ratio conclusion of Theorem~\ref{thm:kappa}. At $\kappa=1$ it recovers the ordinary full-histogram baseline. Sending $\kappa$ to a very large finite value cannot compensate for $Q\ll\sqrt M$. These are representation constraints, not an optimum at fixed total word length, memory, runtime or expanded circuit depth.

\section{Binary sources: active windows and weighted returns}\label{sec:binary}
The three-symbol waiting construction does not describe the entire alphabet dependence. We now set $K=2$, keeping exactly the same one-block interface. The upper requires $D\ge8,r\ge3$, integer $\kappa\ge1$ and legal dyadic $Q\ge16$. The lower permits any fixed $D,r\ge1$ and real $\kappa,Q\ge1$. Its active segment is Markov, not IID, and the useful stationary mixing scale is different.

\subsection{A binary stationary rare-window witness}
Put
\[
 e_0=\max\{\sqrt\beta,M^{-1}\},\qquad z=\kappa Q.
\]
\begin{proposition}[Binary truncated lower]\label{prop:binarylower}
For $M\ge64$, $0\le\beta\le1/256$ and $z\le M/4$,
\begin{equation}\label{eq:binarylower}
 \Risk_{\kappa,\mathrm{stat}}^{(2)}
 \ge\frac{c_{\mathrm{bin}}}{z}
 \ln\left(1+\frac18\min\left\{\frac{z}{Me_0},\frac Mz\right\}\right),
\end{equation}
where $c_{\mathrm{bin}}>0$ depends only on $D,r$. This is \unit{BI001 v2}.
\end{proposition}
\begin{proof}
Order the pair states as $W=00,A=01,B=10,C=11$. Let $a=1/(2M)$, $p\in[1/4,3/4]$, and use the actual homogeneous kernel and stationary law
\[
 P=\begin{pmatrix}
 1-a&a&0&0\\
 0&0&1-p&p\\
 e_0&1-e_0&0&0\\
 0&0&1/2&1/2
 \end{pmatrix},\qquad
 \Pi=x(e_0/a,1,1,2p),\quad x=(e_0/a+2+2p)^{-1}.
\]
Direct multiplication gives stationarity and $\Pi_W\ge4/11$. The two-step entries, together with $e_0^2\ge\beta$ and $e_0\ge4\beta$, give the all-state certificate. For example its restrictive entry is
$P^2(B,A)=e_0a\ge\beta x$; the first row's active entries are at least $a/4$, and the other nonexceptional active entries are at least $1/8$. The proof uses the actual $p$-dependent hidden $\Pi$, not a fixed waiting-state start.

Set $J_*=\floor{\min\{(4e_0)^{-1},M^2/(4z^2)\}}\ge4$.
For $4\le j\le J_*$, consider histories with $M+2-j$ initial zeros, then a first $1$, and an active suffix with no pair $00$. On the active states $A,B,C$ use the reference chain
\[
 R=\begin{pmatrix}0&1-p&p\\1&0&0\\0&1/2&1/2\end{pmatrix},
\]
started at $A$. Relative to its $j-1$-transition law, the actual history density is
\[
 \Pi_W(1-a)^{M-j}a(1-e_0)^{N_B}\ge(32M)^{-1},
\]
where $N_B\le j-1$. The events for different $j$ are disjoint. Their probabilities need not be independent of $p$; pointwise domination, integrated against one common uniform prior, replaces the IID-event argument.

Every row of $R^2$ puts mass at least $1/8$ on $B$. Coupling in two-step blocks therefore gives expected differing-state count at most
$2\sum_{b\ge0}(7/8)^b=16$. For any positional sum $F$ of local terms bounded by $\kappa$, two continuations after a common prefix differ in expected $F$ by at most $34\kappa r$. Doob martingale increments and single-transition differentiation then give
\[
 \Var_p F\le(34\kappa r)^2j,\qquad
 |(\E_pF)'|\le34\kappa rj.
\]
These estimates cover overlapping windows and arbitrary positional tables; no independent-coin argument is used. Also $R^3(s,A)\ge1/8$ for every $s$, so the terminal state is $A$ with probability at least $1/8$ for all $j\ge4$.

Freeze all $H_2=2^{2r}$ boundary branches. With
\[
 T_2=34r\sqrt{64DH_2},\quad N_2=\ceil{2T_2+4}+1,\quad
 A_2=H_2(2+68r)N_2^D,
\]
simultaneous Chebyshev gives a good event of probability at least $31/32$, ratio deviation at most $2\kappa T_2\sqrt j/M$, and curve derivative at most $68\kappa rj/M$. The rescaled cutoff makes the grid-tube radius at most $T_2$. A mesh therefore contains all complete good messages in at most $A_2(1+zj/M)$ labels.

Only terminal $A$ has actual next-symbol parameter $p$. The good-and-$A$ event has probability at least $3/32$ for every $p$. Radius $1/(128R_j)$ neighborhoods about $R_j=A_2(1+zj/M)$ decoder outputs cover prior mass at most $1/32$. Their complement within the good-and-$A$ event gives mean loss at least $1/(2048R_j)$. This is a joint-prior calculation, not a uniform posterior after conditioning on the endpoint.

After pathwise domination and summation,
\[
 \Risk_{\kappa,\mathrm{stat}}^{(2)}
 \ge\frac1{65536A_2M}\sum_{j=4}^{J_*}(1+zj/M)^{-1}.
\]
For $u=z/M\le1/4$, the sum is at least
$\ln(1+uJ_*)/(8u)$, including $J_*=4$.
Since $uJ_*\ge\min\{z/(Me_0),M/z\}/8$, this proves
\eqref{eq:binarylower} with $c_{\mathrm{bin}}=(524288A_2)^{-1}$.
\end{proof}

For this topology the necessity of a square-root mixing scale is visible in
$P^2(B,A)\ge\beta\Pi_A$, which is equivalent to
$e_0^2+e_0a(2+2p)\ge\beta$. Taking both exit and launch probabilities of order $M^{-1}$ only allows $\beta=O(M^{-2})$. This identity alone is not a universal binary obstruction; a matching upper needs a separate argument.

\subsection{Stationary weighted occupation and the actual floor decoder}
Let a stationary binary source have
\[
 \Pi_{00}=y,\quad \Pi_{01}=\Pi_{10}=x,\quad \Pi_{11}=w
\]
and exit parameters
$a=\theta_{00}(1)$, $b=\theta_{01}(1)$,
$e=\theta_{10}(0)$, $t_1=\theta_{11}(0)$.
Define analysis weights
$q_{00}=a$, $q_{11}=t_1$, and $q_{01}=q_{10}=1$.
The crossing weights are clipping constants, not assertions that the crossing rows have unit exit probability.
\begin{lemma}[Binary weighted reciprocal occupation]\label{lem:binaryreturn}
For $\beta>0$, $A>0$ and integer $k\ge1$,
\begin{equation}\label{eq:binaryreturn}
 \E\min\{q_{Y_M},A/(J+k)\}
 \le\frac{4A}{M}\left[17+\ln\left(1+\frac{96}{\sqrt\beta A}\right)\right].
\end{equation}
Zero weights contribute zero before any division. This is \unit{BU001 v2}.
\end{lemma}
\begin{proof}
Stationarity gives $ya=xe$ and $wt_1=xb$. If $x=0$, positive $\beta$ permits only a singleton diagonal support; its weight is zero, so the assertion follows directly. Henceforth assume $x>0$. Positive-$\beta$ minorization makes the supported chain irreducible, and its two-step entries and balance now imply $e^2\ge\beta y$ and $b^2\ge\beta w$. First-step equations for the largest supported hitting mean $H_s$ give, for example,
\[
 H_{00}\le\frac{b}{et_1}+\frac2e+\frac1{t_1}+1,\qquad
 H_{01}\le\frac2a+\frac1{t_1}+2.
\]
Multiplication by $\sqrt\beta\,\Pi_s q_s$, balance and the two square inequalities show
\[
 H_s\le\frac5{\sqrt\beta\,\Pi_s q_s}.
\]
The other targets follow by exchanging symbols. Absent states are omitted using the smaller supported first-step system, not evaluated by $0/0$.

The stationary reverse pair chain becomes the forward chain after swapping pair coordinates. Thus the same bound applies when conditioning on the actual terminal state $s$ and looking backward. Kac gives the mean positive return $1/\Pi_s$. For $\overline H=\max\{H_s,1/\Pi_s\}$,
$\overline H\le6/(\sqrt\beta\,\Pi_sq_s)$.
The elementary hitting-time moment induction yields
\[
 \E\tau^m\le m!\overline H^{m-1}\E\tau .
\]
Indeed $T^m\le m\sum_{j<T}(T-j)^{m-1}$, and on survival the remaining law is a supported hitting law. With
$\lambda=\sqrt\beta\,\Pi_sq_s/24$, this gives
$\E e^{\lambda\tau}\le e^{2\lambda/\Pi_s}$.
Strong Markov at returns and a separate small-$M\Pi_s$ cap imply
\[
 \Prob_s(J<M\Pi_s/16)
 \le2e^{-\sqrt\beta\,M\Pi_sq_s/48}.
\]
This is a weighted return estimate, not a claim of a global $\sqrt\beta$ spectral gap.

Put $\mu=M\Pi_s$ and $\alpha=A/\mu$. The inherited Kac--Markov count bound is
$\Prob_s(J<j)\le j/\mu$. In layer cake for
$\min\{q_s,A/(J+k)\}$, layers below $16\alpha$ cost at most $16\alpha$; higher layers are bounded by
$\min\{\alpha/u,2e^{-\sqrt\beta\mu q_s/48}\}$.
Using $v e^{-\sqrt\beta v/48}\le48/\sqrt\beta$ gives the conditional bound
\[
 \frac{A}{M\Pi_s}
 \left[17+\ln\left(1+\frac{96}{\sqrt\beta A}\right)\right].
\]
Only now multiply by $\Pi_s$ and sum over at most four supported states.
\end{proof}

\begin{proposition}[Unchanged binary program, stationary upper]\label{prop:binaryupper}
For the same paid eight-coordinate triple bank and decoder, every legal $\kappa,Q,M$ and $\beta\in[0,1]$ satisfy
\begin{equation}\label{eq:binaryupper}
 \operatorname{risk}_{\mathrm{stat}}(T_\kappa)
 \le\min\left\{1,\frac{2^{17}}{\sqrt M}+\frac3Q+
 544\,\mathcal W(\kappa Q,Me_0)\right\}.
\end{equation}
This is \unit{BU001 v2}; its proof here is stationary. Section~\ref{sec:initialization} establishes the arbitrary-start extension by a different path/count argument.
\end{proposition}
\begin{proof}
Let $S_0$ be the inherited actual-terminal add-one risk and
$B_0=\E[2/(J+2)]$. Then
\[
 S_0\le2^{13}\sqrt{32}/\sqrt{M+1},\qquad
 B_0\le2\sqrt{32}/\sqrt{M+1}.
\]
For $J>0$, the empirical exit fraction $\widehat p$ differs from the add-one fraction by at most $2/(J+2)$. Thus its expected error from the actual exit probability is at most $S_0+B_0$.

Let $\rho=\kappa J/[M+(\kappa-1)J]$ be the selected mass and $n_0,n_1$ its two integer floors. On $N=n_0+n_1\ge1$, write
$n_0=Q\rho(1-\widehat p)-f_0$,
$n_1=Q\rho\widehat p-f_1$, with $f_i\in[0,1)$.
The normalized exit $v=n_1/N$ obeys
\[
 v-\widehat p=
 \frac{\widehat p f_0-(1-\widehat p)f_1}{N},\qquad
 |v-\widehat p|
 \le\min\{\widehat p,1-\widehat p,2/(Q\rho)\}.
\]
On diagonal rows, clipping at the true exit probability and its Lipschitz dependence charge one empirical error. On crossing rows use the direct bound with $q_s=1$. Since
$2/(Q\rho)\le2/Q+6M/[z(J+2)]$ for $J\ge1$, both cases give
\[
 |v-\theta_{\mathrm{exit}}|
 \le\min\{q_s,A/(J+2)\}+2/Q
       +2|\widehat p-\theta_{\mathrm{exit}}|,\qquad A=6M/z.
\]
The exact weights and true exits here are analysis variables, not added ports.

If the entire selected block floors to zero, $Q\rho<2$ and hence $J<4M/z$; this also includes $J=0$. The unchanged uniform fallback is charged by universal occupation:
\[
 \Prob(N=0)\le16\ceil{4M/z}/(M+1)\le64/z+16/M.
\]
Final ordered output rounding costs at most $1/Q$. Therefore the total risk is at most
\[
 2S_0+2B_0+\E\min\{q_{Y_M},A/(J+2)\}
       +64/z+16/M+3/Q.
\]
For $\beta\ge M^{-2}$, substitute $A=6M/z$ in
\eqref{eq:binaryreturn}; the weighted term is at most
$480[1+\ln(1+z/(M\sqrt\beta))]/z$.
Together with $64/z$ and the statistical bounds this is
\eqref{eq:binaryupper}. For $\beta<M^{-2}$, including zero, use the unchanged beta-free SB001 bound and
$\Psi(96/z)\le96[1+\ln(1+z)]/z$ instead. No positive-$\beta$ return argument is used at zero.
\end{proof}

\subsection{The complete stationary binary frontier and alphabet contrast}
Define
\[
 \mathcal F_2(M,\beta,\kappa,Q)
 =\max\{M^{-1/2},Q^{-1},\mathcal W(\kappa Q,Me_0)\},
 \qquad e_0=\max\{\sqrt\beta,M^{-1}\}.
\]
\begin{theorem}[Stationary binary frontier]\label{thm:binaryjoint}
Fix $K=2$, $D\ge8$ and $r\ge3$. For every $M\ge64$, $\beta\in[0,1]$, integer $\kappa\ge1$ and legal dyadic $Q\ge16$,
\begin{equation}\label{eq:binaryjoint}
 c_2\mathcal F_2\le\Risk_{\kappa,\mathrm{stat}}^{(2)}
 \le\operatorname{risk}_{\mathrm{stat}}(T_\kappa)
 \le\min\{1,131619\,\mathcal F_2\},
\end{equation}
where $c_2=\min\{(256\cdot5)^{-1},c_I/2,c_{\mathrm{bin}}/48\}$ depends only on $D,r$. This is \unit{BS001 v1}. It is a stationary theorem only.
\end{theorem}
\begin{proof}
The uniform IID and statistical floors are already binary witnesses. For the logarithmic term repeat the exhaustive completion in Theorem~\ref{thm:joint}, now with $t=Me_0$ and Proposition~\ref{prop:binarylower}. Specifically, when $\beta>1/256$, $e_0>1/16$ and
$\mathcal W\le Q^{-1}+16/M\le Q^{-1}+2M^{-1/2}$.
The large-$z$ cases $z>M/4$ and $z>M^{3/4}$ depend only on $t\ge1$ and give the same $5M^{-1/2}$ bound. On the complement the binary lower is within its stated domain; $x=z/t\le z\le(M/z)^3$ gives
\[
 \mathcal W(z,t)\le Q^{-1}
 +24z^{-1}\ln(1+\min\{z/t,M/z\}/8).
\]
Thus the three separate lower witnesses imply $c_2\mathcal F_2$ with no uncovered boundary case. The upper follows from \eqref{eq:binaryupper}, since
$2^{17}+3+544=131619$. The same eight-coordinate program supplies admissibility; none of its costs or ports changes.
\end{proof}

The exact stationary statistical-order conditions are consequently
\begin{equation}\label{eq:binarythreshold}
 Q\gtrsim\sqrt M,\qquad
 \kappa Q\gtrsim Z_{*,2},\quad
 Z_{*,2}=\sqrt M\left[1+\ln\left(1+\frac1{\sqrt M e_0}\right)\right].
\end{equation}
The inversion and constructive allocation proof of Corollary~\ref{cor:jointthreshold} applies verbatim with $t=Me_0$: necessity holds for every fixed target multiplier, and sufficiency gives an explicitly large one. This transfers an elementary scalar argument, not the three-symbol lower.

For clarity about the least budget, in a fixed legal band
$c_-\sqrt M\le Q\le c_+\sqrt M$ define
\[
 \kappa_{\min,\mathrm{stat}}^{(2),\mathrm{int}}
 =\min\{k\in\mathbb N:k\ge1,\
       \Risk_{k,\mathrm{stat}}^{(2)}\le C_0/\sqrt M\},
\]
with value $+\infty$ if the set is empty, and put
$f_{Q,2}=1+\ln(1+Q/(Me_0))$.
The lower proof holds also for every finite real envelope. Since
$\mathcal W(\kappa Q,Me_0)\ge f_{Q,2}/(\kappa Q)$, every feasible real budget satisfies
\[
 \kappa\ge\frac{c_2}{C_0c_+}f_{Q,2}.
\]
For the actual integer choice $\kappa=\ceil{f_{Q,2}}\le2f_{Q,2}$,
$\mathcal W(\kappa Q,Me_0)\le2/Q$. Hence the unchanged binary program has risk at most
\[
 \frac{2^{17}+1091/c_-}{\sqrt M}.
\]
For $C_0$ at least this constant, the feasible integer set is nonempty, its least element exists, and
$\kappa_{\min,\mathrm{stat}}^{(2),\mathrm{int}}=\Theta(f_{Q,2})$.
This also bounds the real-budget infimum between the same lower and integer feasible upper; it does not assert a fractional-weight implementation.

\begin{corollary}[A stationary alphabet distinction]\label{cor:alphabet}
At $\beta=1/M$ and a fixed square-root grid band, the least feasible weight ratio is $\Theta(1)$ for $K=2$, but $\Theta(\log M)$ for every fixed $K\ge3$, at sufficiently large fixed target-risk multipliers. Both statements use the same all-state certificate, actual-terminal stationary loss and one-summary-block interface.
\end{corollary}
\begin{proof}
For $K=2$, $e_0=M^{-1/2}$, so $f_{Q,2}=1+\ln(1+Q/\sqrt M)$ is bounded by band-dependent constants. For $K\ge3$, $M\betaz=1$ and $1+\ln(1+Q)=\Theta(\log M)$; Theorems~\ref{thm:kappa} and \ref{thm:joint} supply the comparison. For a chosen fixed $K\ge3$, one may use a common fixed budget $D\ge K^3$, $r\ge3$ for both alphabets. Constants still depend on that architecture; no dimension-uniform or universal Transformer claim follows.
\end{proof}

At $\beta=0$, both alphabets instead require logarithmic weight ratio in this grid band; the binary witness gives the same weak-mixing behavior throughout $\beta\le M^{-2}$. Thus the distinction is not that binary sources avoid rare contexts altogether. The stationary proof above by itself does not establish arbitrary-start refinement: its reversal step uses stationarity, and the older beta-free upper is insufficient. The next section closes that gap without changing the implemented decoder, and extends the alphabet distinction to both initializations.


\section{Initialization through counts, not endpoint mixing}\label{sec:initialization}
The stationary binary proof conditions on the terminal state and reverses a stationary path. Neither operation is available under an arbitrary initial law. A direct replacement by rapid weighted endpoint mixing is false, even with full stationary support. The successful argument instead retains the first hitting time, the actual terminal local time and the source's true row-exit weights.

Throughout this section write $W=00,A=01,B=10,C=11$ and use row parameters
\[
 a=\theta_W(1),\quad b=\theta_A(1),\quad
 e=\theta_B(0),\quad t_1=\theta_C(0).
\]
For arbitrary starts the appropriate analysis weights are
\[
 q=(a,b,e,t_1),
\]
the probabilities of a designated actual successor on each row. In particular the crossing weights are now $b,e$, not the unit caps used in the stationary return lemma. All these quantities are proof variables, never decoder inputs.

\subsection{A negative control for endpoint-only reasoning}
\begin{proposition}[Failure of weighted endpoint transfer]\label{prop:mixingfailure}
There are full-support binary sources with genuine all-state coefficient $\beta_*>0$ for which neither
\[
 \sum_s q_s|(\mu P^n)_s-\Pi_s|,
 \qquad
 \sum_s|(\mu P^n)_s-\Pi_s|(P^nq)(s)
\]
admits a uniform bound
$C n^{-\alpha}+C\exp(-cn\sqrt{\beta_*})$
for any fixed $\alpha,C,c>0$, even when $q$ is the row-minority vector. This is the method obstruction \unit{WM001 v2}, not a precision-risk lower bound.
\end{proposition}
\begin{proof}
Take $0<\eta\le\varepsilon\le1/32$ and
\[
 a=\varepsilon,\quad b=\eta,\quad e=1/2,\quad t_1=\varepsilon\eta,
 \qquad
 \Pi=\frac{(1/2,\varepsilon,\varepsilon,1)}{3/2+2\varepsilon}.
\]
Multiplying the binary pair kernel shows that the exact largest all-state coefficient is
$\beta_*=\eta^2(3/2+2\varepsilon)$; the binding entry is $P^2(A,B)/\Pi_B$. Every other entry dominates it because $\eta\le\varepsilon$. The minority vector is $(\varepsilon,\eta,1/2,\varepsilon\eta)$.

Start at $W$ and reflect an $A\to C$ transition back to $B$. In the reflected $W,A,B$ chain, both crossing-state probabilities are at most $2\varepsilon$ at every time: the recursion is
$r_A(k)=\varepsilon r_W(k-1)+r_A(k-2)/2$.
A first-separation coupling to the original chain consequently gives
\[
 P^n(W,W)\ge1-4\varepsilon-2n\varepsilon\eta .
\]
At $n=\floor{1/(16\varepsilon\eta)}$, this is at least $3/4$, whereas $\Pi_W\le1/3$. Keeping only state $W$ bounds the two displayed residuals below by $5\varepsilon/12$ and $5\varepsilon/16$, respectively. For any fixed $\alpha>0$, choose fixed $k\ge1$ with $\alpha(k+1)>1$ and let $\eta=\varepsilon^k$. Then
$n\asymp\varepsilon^{-k-1}$ and $n\sqrt{\beta_*}\asymp\varepsilon^{-1}$, so the proposed right side is $o(\varepsilon)$.
\end{proof}

Long initial residence can itself supply many observations of the likely terminal row. Both residuals discard this information. Their failure therefore does not separate stationary and arbitrary-start prediction risks; it explains why the following proof must preserve counts.

\subsection{A same-target reciprocal identity}
\begin{lemma}[Schur return identity]\label{lem:schur}
For any homogeneous chain on $d$ states, any initial/target state $s$ and $n\ge1$, let $J_s(n)=\sum_{j<n}\ind\{Y_j=s\}$. Then
\begin{equation}\label{eq:schur}
 \E_s\frac{\ind\{Y_n=s\}}{J_s(n)}
 =\frac{\operatorname{tr}(P^n)-\operatorname{tr}(P_{-s}^n)}n
 \le\frac dn .
\end{equation}
The fraction is evaluated only on $Y_n=s$, where its denominator is positive. No irreducibility, stationarity or aperiodicity is assumed.
\end{lemma}
\begin{proof}
Let $F_s(z)$ be the first-positive-return probability generating function, allowing a defective law. First-step decomposition and Schur complementation give the formal identities
\[
 F_s(z)=zP_{ss}+z^2P_{s,-s}(I-zP_{-s})^{-1}P_{-s,s},
 \qquad
 1-F_s(z)=\frac{\det(I-zP)}{\det(I-zP_{-s})}.
\]
Both determinants have constant term one. Expand their formal logarithms:
\[
 [z^n]\{-\ln(1-F_s(z))\}
 =\frac{\operatorname{tr}(P^n)-\operatorname{tr}(P_{-s}^n)}n.
\]
On the other hand, $-\ln(1-F_s)=\sum_{k\ge1}F_s^k/k$. Strong Markov makes its coefficient
$\sum_k\Prob_s(\tau_k=n)/k$. At the $k$th positive return there are exactly $k$ earlier visits, including time zero. This is the left side of \eqref{eq:schur}. Nonnegativity follows from the return expansion, and the upper bound follows from $\operatorname{tr}(P^n)\le d$ and nonnegativity of the killed trace.
\end{proof}

This identity, proved in \unit{BAI004 v5} and \unit{BOD001}, controls a return to the initial target. For a different terminal target it is used after a first-hit decomposition, rather than as a bound on an independently chosen row.

\subsection{Absorbing diagonals and alternating geometric runs}
Suppose $W$ is absorbing, so $a=0$; no restrictions are placed on the other rows or initial distribution. A true path ending in $A,B,C$ never visited $W$. Its complete-path probability is at most that under the reflected chain
\[
 R=\begin{pmatrix}0&1-b&b\\1&0&0\\0&t_1&1-t_1\end{pmatrix}
\]
on $A,B,C$, since each departure from $B$ contributes an additional factor $1-e\le1$ in the true law. This domination preserves the terminal state and its entire prior count. It is not a conditional IID statement.

For $0<b,t_1<1$, an $A$-run has $G_b$ visits to $A$ and duration $2G_b-1$, while a $C$-run has $G_{t_1}$ visits and duration $G_{t_1}+1$, where the geometric variables are supported on positive integers. Their generating functions and the cycle function are
\[
 D_A(z)=\frac{bz}{1-(1-b)z^2},\quad
 D_C(z)=\frac{t_1z^2}{1-(1-t_1)z},\quad
 F(z)=D_A(z)D_C(z).
\]
Put $L(z)=-\ln(1-F(z))=\sum_{n\ge1}L_nz^n$. Then
\begin{equation}\label{eq:renewalcoeff}
 0\le L_n\le\min\{b,t_1,3/n\}.
\end{equation}
Indeed $L_n=\sum_{k\ge1}\Prob(S_k=n)/k$. Dropping $1/k$ gives a cycle-completion probability; completing at $C$ requires a final $A\to C$ transition, giving the bound $b$. Reversing the order of the independent run durations gives the same cycle law and the bound $t_1$. Finally
\[
 \det(I-zR)=[1-(1-b)z^2][1-(1-t_1)z][1-F(z)]
\]
and the trace/log identity give
$nL_n=\operatorname{tr}(R^n)-(1-t_1)^n
-2\ind_{\{n\ {\rm even}\}}(1-b)^{n/2}\le3$.

Here is the count-sensitive reduction behind the use of these coefficients. For
\[
 g_q(r,N)=\binom{N-1}{r-1}q^r(1-q)^{N-r},\quad N\ge r\ge2,
\]
the exact identity
\[
 \frac{g_q(r,N)}{q(N-1)}
 =\frac{g_q(r-1,N-1)}{r-1}
\]
removes one completed run and one visit. To a path ending at $A$, virtually append $A\to C$, multiplying its probability by $b$; its completed $A$-visit count is $J_A+1$. To a path ending at $C$, append $C\to B\to A$, multiplying by $t_1$, with completed count $J_C+1$. Applying the identity with the larger denominator $1+q(N-1)$ reduces the weighted terminal sums to coefficients of $L$, $D_AL$, $D_CL$ and one-cycle exceptions. The removed durations are exactly one for an $A$-run and two for a $C$-run. No previous visit is reset.

For completeness, splitting the remaining duration at half the horizon gives
\[
 [z^n]F\le8/n,\quad
 [z^M]D_CL\le15/M,\quad [z^M]D_AL\le41/M.
\]
On the short-duration part use $L_{M-\ell}\le6/M$. On the long-duration part use respectively $L_{M-\ell}\le t_1,b$, and geometric density bounds
$t_1e^{-t_1M/4}$ and $be^{-bM/8}$. The resulting maxima
$Mt_1^2e^{-t_1M/4}<9/M$ and $Mb^2e^{-bM/8}<35/M$ give the constants. The run reductions thus imply, for either target $s=A,C$ and every initial active state,
\[
 \E_R\left[\ind\{Y_M=s\}\frac{q_s}{1+q_sJ_s}\right]\le64/M,
 \qquad q_A=b,\quad q_C=t_1,\quad M\ge8.
\]
An initial $B$ shifts the horizon by one; use that shift only for $M\ge9$, and the unit integrand bound at $M=8$. Finite-horizon continuity extends the estimates to $b,t_1\in[0,1]$. Since
\[
 \min\{q,A_0/(j+2)\}
 \le2\max\{A_0,1\}\frac q{1+qj},
\]
the total $A,C$ contribution is at most $256\max\{A_0,1\}/M$.

The true $B$ contribution uses its exit weight $e$. An exit from a terminal $B$ into absorbing $W$ is the first visit to $W$, so
\[
 e\Prob(Y_M=B)
 \le\Prob(J_{Y_{M+1}}(M+1)=0)\le16/(M+2).
\]
The absorbing row has weight zero. Complete-path domination therefore proves
\[
 \E\min\{q_{Y_M},A_0/(J+2)\}
 \le256\max\{A_0,1\}/M+16/(M+2).
\]
Combined with the historywise floor reduction below, this is \unit{BAI003 v2}'s uniform all-start bound
\begin{equation}\label{eq:absorbingrisk}
 \operatorname{risk}(T_\kappa)
 \le\min\{1,2^{17}M^{-1/2}+3/Q+1600/(\kappa Q)\}.
\end{equation}
The earlier \unit{BAI002} reflected-core argument obtains the same no-log order when $b,t_1$ stay in a fixed compact subinterval. The run argument is what removes that extra assumption. Neither argument changes the finite predictor.

\subsection{Full support: diagonal starts and first-hit convolutions}
Assume now $\beta>0$ and full stationary support. All four exits are positive. Write the stationary masses as $(y,x,x,w)$; then
$ya=xe$, $wt_1=xb$ and $x=at_1/(et_1+2at_1+ab)$.
The full all-state certificate supplies
\[
 ae\ge\beta x,\qquad bt_1\ge\beta x,\qquad at_1\ge\beta x.
\]
The third inequality uses $P^2(W,C)=ab\ge\beta w$ and positive $b$.

The diagonal-to-successor-crossing bound in \unit{BAI004 v5} is
\begin{equation}\label{eq:diagcross}
 \E_W\!\left[\ind\{Y_M=A\}\min\{1,A_0/(J_A+2)\}\right]
 \le\frac{40A_0}{M}
 \left[1+\ln\left(1+\frac1{d_0A_0}\right)\right],
 \quad d_0=\max\{\sqrt\beta,M^{-1}\}.
\end{equation}
Its key algebra is the certificate consequence
\[
 \min\{1/e,(1+Mt_1)/b\}
 \le12(1+Ma)/\sqrt\beta,\qquad \beta M^2\ge1.
\]
To see why it is uniform, set
$\alpha=Ma,E=Me,B=Mb,T=Mt_1$ and
$D=\alpha B+2\alpha T+ET$. The three certificate inequalities become
$\beta M^2\le D,ED/T,BD/\alpha$.
For $H=\max\{E,B/(1+T)\}$ and $X=(1+\alpha)H$, splitting at $T=1$ and $\alpha=1$ gives
$\beta M^2\le6X^2+6X$; hence $X\ge M\sqrt\beta/12$.

Before an excursion from $A$ hits $W$, its expected truncated number of $A$ visits is at most
$\min\{1/e,(1+Mt_1)/b\}$.
The first bound counts the geometric tests at $B$; the second counts outgoing $A\to C$ entries and uses that they are at most the number of $C$ exits plus one. Include the outgoing transition after a possible visit at the observation endpoint. In the full chain there are at most $Ma$ further $W\to A$ entrances in expectation. Thus the global visit count $V_A$ obeys
\[
 \E_A V_A\le12(1+Ma)^2/\sqrt\beta .
\]
For every path, summing the terminal functional over all horizons counts successive visits once:
\[
 \sum_{n=0}^M
 \ind\{Y_n=A\}\min\{1,A_0/(J_A(n)+2)\}
 \le A_0[1+\ln(1+V_A/A_0)].
\]
Jensen applies to this \emph{global} count, not to a sum of reset excursion counts.

From initial $W$, the first hit of $A$ is geometric$(a)$. Split its convolution with the same-target functional at $M/2$. Early hits cost at most $8A_0/M$ by Lemma~\ref{lem:schur}. Late hits have point mass at most $ae^{-aM/4}$. Multiplying the preceding cumulative-count bound and using
$x e^{-x/4}\le2$, $x^2e^{-x/4}\le9$
proves \eqref{eq:diagcross}. If $\beta M^2<1$, use the universal occupation bound directly; it gives the same cutoff and avoids every zero divisor.

For the other crossing target, every initial-$W$, terminal-$B$ path has a last $A$ visit, at time $k$, and
$J_B(M)=J_A(k)$. The intervening phase has distribution
\[
 U_C(1)=1-b,\qquad U_C(\ell)=bt_1(1-t_1)^{\ell-2},\quad \ell\ge2.
\]
Consequently its terminal count functional is exactly the convolution of \eqref{eq:diagcross}'s functional at time $k$ with $U_C(M-k)$. For $k\ge M/2$ use \eqref{eq:diagcross}; for $k<M/2$ use the long-hold density
$bt_1e^{-t_1M/4}$ and the global bound
$b\E_WV_A\le1+Mt_1$. The term $b\ln(1/b)$ stays multiplied by $b$ and is bounded by $1/\exp(1)$. This proves the opposite-crossing estimate with constant $100$ in place of $40$. Horizons $8\le M<16$ use the direct cap, and the weak-$\beta$ branch is separated before dividing by $b$. Symbol exchange supplies both crossing targets from initial $C$.

\subsection{Opposite diagonals: a density envelope and a hitting potential}
The remaining component is initial $W$, terminal $C$, with true exit weight $t_1$. Let $T_C$ be the first hit of $C$, $\rho=b+e-be$, and
\[
 m=\E_WT_C=\frac{\rho+a(2-b)}{ab},\qquad \lambda=1/m.
\]
The first-step equations give a largest supported hitting mean $H$ with
$m\le H\le m+1\le2m$. Blockwise Markov and strong Markov therefore yield
$\Prob_i(T_C>n)\le2e^{-\lambda n/16}$.
This survival estimate alone is insufficient; the required first-hit \emph{density} has the stronger envelope
\begin{equation}\label{eq:hitdensity}
 d_n=\Prob_W(T_C=n)\le12\lambda e^{-\lambda n/16}.
\end{equation}
Indeed the killed probabilities satisfy
$A_k=aW_{k-1}+(1-e)(1-b)A_{k-2}$, giving
\[
 d_n=ab\sum_{j=0}^{\floor{(n-2)/2}}(1-\rho)^jW_{n-2-2j}.
\]
Since $\lambda\le b\le\rho$, the survival envelope bounds this by
$4(ab/\rho)e^{-\lambda n/16}$; independently it is at most
$3be^{-\lambda n/16}$. The identity
$\min\{b,ab/\rho\}/\lambda
=\min\{1,a/\rho\}(\rho/a+2-b)\le3$
proves \eqref{eq:hitdensity}.

Let $p=\Pi_C$ and $h_i=\E_iT_C$, $h_C=0$. The hitting potential satisfies
\[
 Ph-h=-1+\ind_{\{C\}}/p.
\]
At $C$ this is Kac's positive-return identity; elsewhere it is the first-step equation. Telescoping from initial $C$ yields, for $V_C=\sum_{n=0}^M\ind\{Y_n=C\}$,
\[
 \E_CV_C=p[M+1+\E_Ch(Y_{M+1})]\le p(M+1+H).
\]
The weighted hitting bound from Lemma~\ref{lem:binaryreturn}'s proof gives
$\sqrt\beta\,pt_1H\le5$, hence
\[
 \E_C(t_1V_C)\le\frac{10}{\sqrt\beta}(1+M\lambda).
\]
The weight cannot be dropped.

Put $f(n)=\E_C[\ind\{Y_n=C\}\min\{t_1,A_0/(J_C(n)+2)\}]$.
Lemma~\ref{lem:schur} gives $f(n)\le4A_0/n$ for $n\ge1$, while chronological visit summation and Jensen give
\[
 \sum_{n=0}^Mf(n)
 \le A_0\left[1+\ln\left(1+
 \frac{10(1+M\lambda)}{\sqrt\beta A_0}\right)\right].
\]
There are no target visits before $T_C$, so the desired quantity equals
$\sum_{r=1}^Md_rf(M-r)$, including $f(0)$ when the first hit is terminal.
Early hits cost $8A_0/M$; late hits use \eqref{eq:hitdensity} and the cumulative bound. With $x=M\lambda$, the maxima
$x e^{-x/32}\le12$ and $x^2e^{-x/32}\le555$
give
\begin{equation}\label{eq:oppositediag}
 \E_W[\ind\{Y_M=C\}\min\{t_1,A_0/(J_C+2)\}]
 \le\frac{8192A_0}{M}
 \left[1+\ln\left(1+\frac1{\sqrt\beta A_0}\right)\right].
\end{equation}
This is \unit{BOD001}; it needs no condition $\beta M^2\ge1$. That fact matters when applying it to a shorter suffix.


\subsection{All initial states, all supports, and the unchanged decoder}
For $\beta>0$, stationarity gives $\Pi=(y,x,x,w)$.
If $x>0$, the all-state entries $ea\ge\beta x$ and $bt_1\ge\beta x$ force all four exits positive, and balance gives $y,w>0$. If $x=0$, supported diagonals are absorbing; two supported absorbing diagonals contradict positive minorization. Thus the stationary support is either full or a single absorbing diagonal. The latter is already covered by \eqref{eq:absorbingrisk}, including initial states outside that support. At $\beta=0$ no support classification is required.

For full support put
\[
 V_M=\sum_s\ind\{Y_M=s\}\min\{q_s,A_0/(J_s(M)+2)\},
 \qquad G=1+\ln(1+1/(\sqrt\beta A_0)).
\]
From either diagonal start, the same-target, successor-crossing, other-crossing and opposite-diagonal bounds have constants $4,40,100,8192$. Hence, at every horizon $m\ge8$,
\[
 \E_{\mathrm{diag}}V_m\le8336 A_0G/m.
\]
The crossing bounds use $\max\{\sqrt\beta,1/m\}$, whose larger value only decreases the logarithm; \eqref{eq:oppositediag} has no strong-certificate horizon restriction. This proves the bound even for the shorter suffixes below.

For a crossing start, let $h_c=b+e-be$. Before the first diagonal hit the chain alternates between $A,B$. If $h_c\le M^{-1/2}$, the crossing endpoint weights contribute at most $b+e$. Last-entry geometric sums give
$a\Prob(Y_M=W)\le e$ and $t_1\Prob(Y_M=C)\le b$; a zero exit contributes zero before any division. Since $b+e\le2h_c$,
$\E V_M\le4M^{-1/2}$.
If $h_c>M^{-1/2}$, the probability of no diagonal hit by $\floor{M/2}$ is at most
$e^{-\sqrt M/8}\le3M^{-1/2}$.
On its complement, strong Markov gives an ordinary suffix starting at a diagonal, of length at least $\ceil{M/2}$. The full terminal count dominates its suffix count. Consequently, for all initial laws,
\begin{equation}\label{eq:allstartoccupation}
 \E V_M\le4M^{-1/2}+\frac{16672A_0}{M}G,
 \qquad \beta M^2\ge1.
\end{equation}
This is a path/count argument, not the false weighted endpoint transfer of Proposition~\ref{prop:mixingfailure}.

The binary floor identity in Proposition~\ref{prop:binaryupper} is historywise. For any chosen true exit probability $q_s$, including $q_A=b,q_B=e$, its positive-block bound is
\[
 |v-\widehat p|\le\min\{\widehat p,1-\widehat p,2/(Q\rho)\}
 \le\min\{q_s,2/(Q\rho)\}+|\widehat p-q_s|.
\]
The arbitrary-start add-one comparator and universal occupation bound therefore give exactly
\begin{equation}\label{eq:allstartdecoder}
 R(T_\kappa)\le
 2S_0+2B_0+\E V_M+64/z+16/M+3/Q,\qquad A_0=6M/z.
\end{equation}
The separate empty-floor charge includes $J=0$ and does not change the uniform fallback. No stationary return estimate is being imported under a nonstationary law.

For full support and $\beta M^2\ge1$, substitution in \eqref{eq:allstartoccupation} gives numerical coefficient
$6\cdot16672+64=100096<2^{17}$
on $\mathcal W(z,Me_0)$, where $e_0=\max\{\sqrt\beta,M^{-1}\}$. The statistical coefficient
$(2^{14}+4)\sqrt{32}+6$ is also below $2^{17}$.
Singleton support uses \eqref{eq:absorbingrisk}. For $0\le\beta<M^{-2}$, use the unchanged beta-free SB001 upper, including $\Psi(96/z)\le96[1+\ln(1+z)]/z$; here $Me_0=1$.
Thus every support, initialization and cutoff boundary is accounted for.

\begin{theorem}[Binary one-block frontier for both initializations]\label{thm:binaryall}
In the family of Theorem~\ref{thm:binaryjoint}, with the same legal budgets,
\begin{equation}\label{eq:binaryall}
 c_2\mathcal F_2\le
 \Risk_{\kappa,\mathrm{stat}}^{(2)}
 \le\Risk_{\kappa,\mathrm{all}}^{(2)}
 \le\operatorname{risk}_{\mathrm{all}}(T_\kappa)
 \le\min\left\{1,\frac{2^{17}}{\sqrt M}+\frac3Q+
 2^{17}\mathcal W(\kappa Q,Me_0)\right\}.
\end{equation}
In particular both minimax risks are of order $\mathcal F_2$, with common upper constant $2^{18}+3=262147$. This is \unit{BIF001 v1}.
\end{theorem}
\begin{proof}
The preceding support and decoder argument proves the upper. Theorem~\ref{thm:binaryjoint}'s stationary lower applies to the same encoder family, and stationary laws are an admissible subset of arbitrary initializations. No new lower witness or all-regime extrapolation is needed.
\end{proof}

Both simultaneous budgets in \eqref{eq:binarythreshold} now characterize statistical-order risk for either initialization. In a fixed square-root grid band, every feasible real weight envelope is at least
$c_2f_{Q,2}/(C_0c_+)$, while the actual integer choice
$\ceil{f_{Q,2}}$ has arbitrary-start risk at most
\[
 \frac{2^{17}+(3+2^{18})/c_-}{\sqrt M}.
\]
Thus the least feasible integer weight ratio is $\Theta(f_{Q,2})$ for sufficiently large fixed target multipliers in both experiments. Corollary~\ref{cor:alphabet}'s constant-versus-logarithmic distinction at $\beta=1/M$ consequently holds also for arbitrary initialization. The previously established stationary constant $544$ remains sharper than the common envelope; it is not invalidated by the extension.

\section{Binary fixed-depth causal finite evaluators}\label{sec:binarycausal}
The one-block real-table family and the causal finite-evaluator family are different. The binary frontier transfers to the latter only after proving a nonlinear Markov concentration bound and an exact finite-label upper inclusion.

Use the temporal-causality and bounded-primitive conditions of Section~\ref{sec:experiment}, now with $K=2$, fixed radius $r\ge3$ and output budget $D\ge8$. The capacity must admit one global head, width $16$, supplied exact radius-three triple-bit tables, and affine coefficient and row-sum budgets at least two. The main prototype uses row-sum-one projections. Computing the triple bits instead as $\operatorname{ReLU}(u+v+w-2)$ is an optional different local implementation requiring row-sum budget three. Public masks must allow density $c\le1/3$. Historical tables and states cannot use the future boundary label; that label enters only the appended query and decoder. Evaluation remains the fixed-architecture FA001 procedure at target $h+3$.

\begin{theorem}[Binary bounded causal frontier]\label{thm:binarycausal}
Put
\[
 H_2(M,\beta,Q)=
 \max\{M^{-1/2},\mathcal W(Q,M\max\{\sqrt\beta,M^{-1}\})\}.
\]
For this specified fixed-depth causal finite-evaluator family,
\[
 c_{\mathrm{arch},2}H_2
 \le\Risk_{\mathrm{causal},\mathrm{stat}}^{(2)}
 \le\Risk_{\mathrm{causal},\mathrm{all}}^{(2)}
 \le\min\{1,262147\,H_2\}.
\]
The stationary lower and exact finite prototype are \unit{BC001 v1}; the arbitrary-start upper composition is \unit{BIF001 v1}. Constants depend on the fixed architecture, not on $M,\beta,Q$.
\end{theorem}
\begin{proof}
For the lower, retain the binary source and active reference chain of Proposition~\ref{prop:binarylower}. Section~\ref{sec:causal} bounds the ideal terminal output's change under any late single-symbol replacement by $\Gamma/M$, already including local preprocessing. The reference chain's two-step coupling has expected disagreement count at most $16$. Telescoping single-symbol replacements along two coupled continuations therefore bounds the difference of their conditional output expectations by $34\Gamma/M$. For a $j$-step active suffix and each frozen branch, the Doob and single-transition derivative arguments give
\[
 \Var_\nu f\le(34\Gamma)^2j/M^2,\qquad
 |(\E_\nu f)'|\le34\Gamma j/M.
\]
This concerns the expectation of the nonlinear ideal output, not a ratio of expectations or a derivative of finite rounding.

Write $H_b=2^b$ for the finite boundary alphabet and
$\tau_B=34\Gamma\sqrt{32DH_b}$. On a simultaneous event of probability at least $31/32$, deviations are at most $\tau_B\sqrt j/M$. With
$J_*=\floor{\min\{(4e_0)^{-1},M^2/(4Q^2)\}}$,
the scaled tube has bounded cross-section. A mesh gives at most
\[
 A_B(1+Qj/M),\qquad
 A_B=H_b(2+34\Gamma)
      \bigl(\ceil{2\tau_B+6}+1\bigr)^D
\]
complete good labels, including the certified finite error at most $1/(8Q)$ and final grid floor. The informative-terminal-$A$ Bayes calculation and pathwise domination from Section~\ref{sec:binary} then yield the same truncated logarithmic lower, with constant $(524288A_B)^{-1}$ and $z=Q$.
The fixed-architecture IID curve argument and full-history testing supply separate $1/Q$ and $M^{-1/2}$ floors. The binary all-regime completion proves $c_{\mathrm{arch},2}H_2$. No growing-weight uniform IID strengthening is needed for this fixed architecture.

For upper membership, supply the eight exact local triple bits at each historical row. Set the synthetic initial state to zero. A single head has zero query and key projections, values those bits, and masks
\[
 \mathcal K_1=\{1\},\quad \mathcal K_2=\{1,2\},\quad
 \mathcal K_t=\{3,\ldots,t\}\ (3\le t\le M+2),\quad
 \mathcal K_{M+3}=\{3,\ldots,M+2\}.
\]
These masks are public, causal, boundary-independent and have density at least $1/3$. Width $16$ holds the old eight feature slots and eight output slots. Projections and identity padding are exact; historical states do not use the terminal pair.

Crucially, the unchanged shifted Taylor evaluator assigns the \emph{same positive integer} $W_*$ to every zero-logit edge. It need not equal one. For fine denominator $D_f=2^P$ and count $C_{abc}$, the exact accumulated denominator and numerator are
\[
 Z=MW_*,\qquad N_{abc}=W_*D_fC_{abc},\qquad
 O_{abc}=\floor{D_fC_{abc}/M}.
\]
Since $D_f/Q$ is an integer,
\[
 \floor{O_{abc}/(D_f/Q)}=\floor{QC_{abc}/M}.
\]
Thus the final bank is exactly the uniform histogram bank at $\kappa=1$ on every history, not merely close to it. Zero counts remain zero, which is necessary for the binary purity argument. The boundary is used only for post-grid block selection. Theorem~\ref{thm:binaryall}'s upper now applies at $z=Q$ for every initialization, although the encoder implementations are different. Since $\mathcal W(Q,Me_0)\ge1/Q$, the claimed envelope follows.
\end{proof}

The finite-cost composition also remains distinct. Charge the complete FA encoder, local-table and mask generation, all-token arrays, every Taylor/edge scan, compact converted records and full arithmetic words from Section~\ref{sec:costs}. Add bank conversion and all-token pointwise decoding. The complete addressed SB001 ledger at $\kappa=1$ is a conservative overcharge for its decoder suffix; using that bound does \emph{not} execute singleton heads inside the dense causal encoder. Add simultaneous phase costs, rather than taking minima across implementations. Optional PURE unrolling applies to the composed bitwork and has much larger depth.

At fixed depth the necessary and sufficient binary grid order is $Z_{*,2}$ in \eqref{eq:binarythreshold}, for both initializations. In particular $\beta=1/M$ permits $Q\asymp\sqrt M$, whereas $\beta\le M^{-2}$ requires $Q\asymp\sqrt M\log M$. The larger-alphabet bounded causal theorem retains the $\beta_0$ cutoff. These distinctions do not depend on an online estimate of the mixing promise.

\section{Actual bounded-primitive depth as a third resource}\label{sec:depth}
Increasing a pre-cut representation's sensitivity need not require growing attention weights. Repeated bounded affine/ReLU transformations can redistribute grid cells toward small histogram coordinates. The resulting depth question is different from the one-block weight/grid frontier: it concerns actual sequential primitive stages before the same information cut. The construction below uses only zero attention logits, so its attention-weight ratio is one.

Scalar companding---a monotone transform, uniform quantization and inverse reconstruction---is established. Adler, Tang and Polyanskiy analyze probability-vector companders under KL loss, including power transforms and the square-root case \citep[Sections I--III, Theorem 5]{adler2022}. The obligations here are instead actual-terminal Markov TV risk, a bounded-primitive dyadic realization, directional finite arithmetic and a matching depth-order lower. Likewise, depth versus \emph{coefficient} precision is an established ReLU-network tradeoff \citep[Proposition 4.1]{ou2024}; our final-summary grid is a different resource, and internal coefficient precision and memory remain paid.

\subsection{Uniform family and theorem}
Fix the alphabet, width, head count, output dimension $D$, local radius $r$, number $b_0$ of boundary-symbol slots, fan-ins, coefficient/operator bounds, state ranges, mask density and number of primitive affine/ReLU operations per stage. Only the number $d\ge1$ of stages may grow. Attention-free positionwise stages are allowed; a syntax requiring a head at every stage may instead execute paid zero-output heads. A stage may not contain an uncharged growing circuit.

The initial synthetic terminal state is a function \emph{only} of public data and the frozen finite boundary field $S$. Every source symbol used by that initial query, including any terminal local window, must be included in the $b_0$ slots before branch counting. Historical local tables occur strictly before the synthetic terminal index $T=M+3$. Every frozen branch obeys the primitive and ideal-range contracts on every formal history, even if its boundary symbols are inconsistent with that history. The cut and decoder access restrictions are unchanged.

For lower bounds, a finite evaluator must have a certified uniform pre-cut error at most $E/Q$, with $E$ independent of $d,M,Q$. We first prove such a certificate for the same signed Taylor/affine evaluator used earlier. Arbitrary discontinuous rounded programs without that certificate are not included.

\begin{theorem}[Depth at the statistical grid]\label{thm:depth}
Within a uniform bounded-primitive family having sufficient fixed capacity for the complete $K^3$ histogram and the coordinatewise construction below:
\begin{enumerate}
\item For every fixed $K\ge2$, every $\beta\in[0,1]$ and every initialization, a fully paid finite encoder with zero attention logits, fixed width and $d=O(\log\log M)$ attains risk $O_K(M^{-1/2})$ at a dyadic $Q\asymp\sqrt M$.
\item For fixed $K\ge3$, $\beta=0$ and a fixed band
$c_-\sqrt M\le Q\le c_+\sqrt M$ with $0<c_-\le c_+<\infty$, attaining any fixed risk bound $C_{\rm risk}/\sqrt M$ requires $d=\Omega(\log\log M)$, even under stationary initialization. This holds for the uniformly certified growing-depth finite evaluator as well as the ideal family.
\end{enumerate}
Thus the least sufficient primitive depth is $\Theta(\log\log M)$ at this operating point and sufficiently large fixed target-risk multiplier. This is not an all-$Q$ three-resource frontier, a necessary depth claim for every positive $\beta$, or a statement about the depth of a Boolean/PURE expansion.
\end{theorem}

\subsection{Dyadic compander geometry}
Use $L\ge0$ for the number of compander iterations and write $\chi=2^L$ for its bottom-segment gain. This symbol is deliberately distinct from the attention-weight ratio $\kappa$. On $[0,1]$ define
\begin{equation}\label{eq:compander}
 T(x)=\min\{1,4x\}=4x-\operatorname{ReLU}(4x-1),\qquad
 g_0(x)=x,\quad
 g_{L+1}(x)=\tfrac12g_L(T(x))+\operatorname{ReLU}(x-\tfrac14).
\end{equation}
A sequential realization sets $u_0=x,v_0=0$ and, for $0\le j<L$, performs
\begin{equation}\label{eq:companderrails}
 u_{j+1}=T(u_j),\qquad
 v_{j+1}=v_j+2^{-j}\operatorname{ReLU}(u_j-\tfrac14),
 \qquad g_L(x)=v_L+2^{-L}u_L.
\end{equation}
The old $u_j$ is retained until both updates finish. Every iteration has a fixed number of bounded-fan-in affine/ReLU primitives. Coefficients have magnitude at most four; $u\in[0,1]$, $v,g_L\in[0,3/2]$, and affine temporaries have a uniform bound, for example four. The dyadics $2^{-j}$ have actual description length $O(j+1)$, which is charged.

Induction gives a strictly increasing concave piecewise-linear function with slope $2^j$ on $[4^{-j-1},4^{-j}]$, $0\le j<L$, and slope $\chi$ on $[0,4^{-L}]$. Its inverse slopes obey $2^{-j}\le2\sqrt x$ outside the bottom segment and equal $1/\chi$ on that segment. Integrating, including across breakpoints, yields
\begin{equation}\label{eq:companderinverse}
 x-y\le(2\sqrt x+\chi^{-1})\,[g_L(x)-g_L(y)],
 \qquad 0\le y\le x\le1.
\end{equation}
This inequality is a proof device; the program never evaluates a square root.

Acquire radius-three one-hot triple features by paid local tables or the paid predecessor-copy construction. A zero-logit head pools the complete $K^3$ histogram at every row over that row's own causal prefix, using the density-$1/3$ masks from Section~\ref{sec:binarycausal}. No terminal counts are copied backward into historical states. Apply $g_L/2$ to every coordinate, then make the downward $Q$-grid cut. The decoder receives only this bank, the exact terminal pair and public metadata. It inverts $g_L$ on twice the represented coordinates, selects the terminal pair's block, normalizes if nonzero, returns uniform otherwise, and performs ordered simplex rounding.

\subsection{Downward arithmetic and zero purity}
A particularly transparent implementation uses
\[
 P=h+L+4,\qquad D_f=2^P=16\chi Q.
\]
For an exact histogram count $A$ and row length $m$, put
$X=\floor{D_fA/m}$, so $\widehat x=X/D_f\le x=A/m$.
Identical positive finite softmax weights cancel exactly, just as in Section~\ref{sec:binarycausal}. In integers, initialize $U_0=X,V_0=0$ and compute
\begin{align}\label{eq:downrails}
 R_j&=\max\{0,U_j-D_f/4\},&
 V_{j+1}&=V_j+\floor{R_j/2^j},&
 U_{j+1}&=\min\{D_f,4U_j\},\\
 W&=V_L+\floor{U_L/2^L}.&&&&\notag
\end{align}
All $T$ updates are exact on the fine lattice. Each positive dyadic-term floor loses less than $1/D_f$. Since $g_L$ is $\chi$-Lipschitz,
\begin{equation}\label{eq:downerror}
 0\le g_L(x)-W/D_f
 \le\frac{\chi+L+1}{D_f}\le\frac{3\chi}{D_f}.
\end{equation}
For $x=0$, every rail and output is exactly zero: there is no positive-offset leakage. Emit $Z=\floor{QW/(2D_f)}$. Then
\[
 0\le g_L(x)-2Z/Q<3/Q.
\]
Consequently $y=g_L^{-1}(2Z/Q)$ is defined and lies below $x$, with
\begin{equation}\label{eq:inverseerror}
 0\le x-y\le\frac{6\sqrt x+3/\chi}{Q}.
\end{equation}
The one-sided inequality, not merely a symmetric error estimate, is what makes the zero-bank and small-mass analysis simple.

\subsection{A uniform extension of the finite evaluator}
The preceding minimal-$P$ schedule is useful, but inclusion in the comparison family must also hold for the same generic evaluator used by the lower. One cannot substitute a growing $d$ into the old fixed constant $C_{\rm arch}$.

Fixed \emph{per-stage} bounds in the affine/ReLU/Taylor argument give computable rational constants $\alpha,\zeta\ge0$, independent of $d,M,h$, such that the numerical discrepancy satisfies
\[
 \delta_{i+1}\le\alpha\delta_i+\zeta(P+1)2^{-P}
 \quad\hbox{whenever }\delta_i\le1,\qquad
 \delta_0\le2^{-P}.
\]
The score neighborhood, positive shift and Taylor index offset are uniform because the primitive bounds are uniform. Enlarge these constants to cover the fixed readout as one additional stage. Choose integer $A_*\ge\max\{2,\alpha+\zeta+1\}$ and then one integer $c_*\ge4$ satisfying
$A_*(c_*+5)2^{1-c_*}\le1$; finite search finds it. With
\begin{equation}\label{eq:depthprecision}
 t=d+2,\qquad P=4h+c_*t,
\end{equation}
the scalar recurrence gives an error bound $A_*^t(P+1)2^{-P}$. Indeed,
\[
 A_*^t(P+1)2^{-P+h}
 \le A_*^t(c_*t+5)2^{-c_*t}
 \le(c_*+5)\,[A_*2^{1-c_*}]^t\le1.
\]
Here $(4h+c_*t+1)2^{-3h}\le c_*t+5$ and
$c_*t+5\le(c_*+5)2^t$. Applying the same estimate to each earlier stage gives $\delta_i\le2^{-h}\le1$, closing rather than assuming the neighborhood bootstrap. The pre-cut error is therefore at most $1/Q$, uniformly in depth and every formal history.

For the compander choose the maximum of \eqref{eq:depthprecision} and $h+L+4$; increasing $P$ preserves the estimate. Exact local one-hot inputs and zero logits again yield the exact downward fine histogram. On positive rails, $T$ remains exact; only the designated dyadic products round downward. A separately floored final division by two adds at most $1/D_f$ in $g_L/2$ space, so the generic evaluator obeys
\[
 0\le g_L(x)-2F_{\rm fin}(x)
 \le\frac{\chi+L+3}{D_f}\le\frac{5\chi}{D_f}.
\]
After the final $Q$ floor, the total loss is less than $(2+5/16)/Q<3/Q$. Hence \eqref{eq:inverseerror} and exact zero purity hold for this \emph{same} uniformly certified evaluator. Unused zero-value heads produce zero exactly. This is a sufficient $P=4h+O(d+1)$ certificate, not a minimal working-precision theorem.

\subsection{Normalization and arbitrary-start occupation}
Let the selected true histogram coordinates be $x_c=C_c/M$, whose mass is $r=J/M$. For $J\ge1$, the removed mass $e=\sum_c(x_c-y_c)$ satisfies
\[
 e\le6\sqrt{Kr}/Q+3K/(\chi Q).
\]
If $e<r$, the two-term normalization triangle inequality gives
$\TV(y/\sum y,x/r)\le e/r$. If $e=r$, the decoded bank is zero and any simplex fallback has distance at most $1=e/r$. Thus, including this branch,
\begin{align}\label{eq:compandernorm}
 \TV(\widehat p_{\rm raw},C/J)
 &\le\frac{6\sqrt{K(K+1)M}}{Q\sqrt{J+K}}
   +\min\left\{1,\frac{3K(K+1)M}{\chi Q(J+K)}\right\}.
\end{align}
Adding the add-one difference $K/(J+K)$ and final rounding $(K-1)/Q$ also covers $J=0$.

The universal occupation bound
$\Prob(J<w)\le\min\{1,K^4w/(M+1)\}$ for integer $w\ge1$
already holds for every initialization. Besides the reciprocal bound expressed by $\Psi$, it gives
\begin{equation}\label{eq:fractionaloccupation}
 \E(J+K)^{-1/2}\le\frac{2K^2}{\sqrt{M+1}}.
\end{equation}
To check the fractional form, for $u>K$ the event $J+K<u$ equals
$J<\ceil{u-K}$, whose positive integer threshold is at most $u$.
For $u\le K$ it is empty. With $a_0=K^4/(M+1)$, layer cake therefore gives
\[
 \E(J+K)^{-1/2}
 \le\int_0^\infty\min\{1,a_0/v^2\}\,dv=2\sqrt{a_0}.
\]
Using \eqref{eq:compandernorm}, \eqref{eq:fractionaloccupation} and the inherited ORIGINAL add-one risk yields
\begin{equation}\label{eq:depthupper}
 \begin{split}
 \Risk\le\min\bigg\{1,\;&2^{13}H_{\rm stat}
 +\Psi\left(\frac{K^5}{M+1}\right)
 +\frac{12K^2\sqrt{K(K+1)}+K-1}{Q}\\
 &+\Psi\left(\frac{3K^5(K+1)}{\chi Q}\right)\bigg\},
 \qquad H_{\rm stat}\le\frac{K^{5/2}}{\sqrt{M+1}}.
 \end{split}
\end{equation}
Here $\Psi(x)=x(1+\ln(1/x))$ on $0<x<1$, $\Psi(0)=0$ and $\Psi(x)=1$ for $x\ge1$, as earlier. This step inherits the same statistical input as SB001; it is not an independent new proof of that input.

Set
\[
 m=\bl(M+1),\qquad h=\max\{4,\ceil{m/2}\},\qquad
 L=\ceil{\log_2m}=\bl(m-1).
\]
These choices require only integer bitlengths. Then
$\sqrt{M+1}\le Q\le2\sqrt{M+1}$ and $m\le\chi<2m$.
The elementary $\Psi$ estimates used for the triple bank imply
$\Psi(C/(mQ))\le3C/Q$ and
$\Psi(K^5/(M+1))\le2K^{5/2}/\sqrt{M+1}$.
Equation~\eqref{eq:depthupper} is therefore $O_K(M^{-1/2})$.
The complete primitive depth is $d=O(L+1)$, counting every affine/ReLU composition. Neither $\beta$ nor a transcendental oracle is used.

\subsection{A depth-independent rescaled lower tube}
Let the uniform per-stage bounds be those of Section~\ref{sec:causal}: at most $H$ heads, values bounded by $V$, score and value Lipschitz bounds $S_*,V_L$, row-update Lipschitz bound $F_L$, local-table diameter $\Delta$ and terminal projection bound $G$. Put
\[
 C_{\rm att}=e^{2B}/c_{\rm mask},\quad
 a=F_L(1+2HVC_{\rm att}S_*),\quad
 b=F_LHC_{\rm att}(V_L+2VS_*),\quad w_0=1+\ln2.
\]
Increase $a$ if needed to cover attention-free stages. The nonnegative matrix recurrence from Section~\ref{sec:causal}, applied only to late source changes, gives
\begin{equation}\label{eq:depthgain}
 \|\Delta_iF\|_\infty\le\gamma_d/M,\qquad
 \gamma_d=\max\{1,rG\Delta[(a+b)^d-a^d]w_0^{d-1}\}\le C_*^d
\end{equation}
for $i\ge T/2$, where the fixed constant
$C_*=\max\{2,\max(1,rG\Delta)\max(1,(a+b)w_0)\}$ suffices.
Because $S$ is frozen, the changed source coin affects at most $r$ historical initial states, all strictly before $T$, and does not change the synthetic initial state. Thus the estimate uses off-diagonal entries only; no diagonal $a^d$ terminal-input contribution has been silently discarded.

Write $\gamma=\gamma_d$, $z=\gamma Q$, $H_b=3^{b_0}$ and $\tau=\sqrt{5DH_b}$. On an IID active suffix of length $j$ in the late half, each frozen branch and coordinate satisfies
\[
 |(\E_pF)'|\le\gamma j/M,\qquad
 \Var_p(F)\le\gamma^2j/(2M^2).
\]
A union bound gives simultaneous deviations at most
$\tau\gamma\sqrt j/M$ with probability at least $9/10$.
If $z\sqrt j/M\le1/2$, a parameter mesh with spacing $M/(zj)$ and the certified $E/Q$ arithmetic error give at most
\begin{equation}\label{eq:depthlabels}
 A_{\rm box}(1+zj/M),\qquad
 A_{\rm box}=2H_bN_{\rm box}^D,\quad
 N_{\rm box}=\ceil{2\tau+2E+4}+1
\end{equation}
complete good messages. The actual boundary label is selected after sampling the unconditional product suffix. All constants in \eqref{eq:depthlabels} are independent of $d$ and $\gamma$. In particular the old depth-dependent global square-sum integral is not being absorbed into a constant.

For $0\le\beta\le1/16$, use the three-symbol stationary witness of Section~\ref{sec:lower}. When $1\le z\le M/4$, restrict to
\[
 2\le j\le J_*=\floor{\min\{(4\beta_0)^{-1},M^2/(4z^2)\}}.
\]
Every such suffix is late, has probability at least $1/(32M)$ independently of the common prior parameter, and satisfies the tube condition. The complete-message Bayes lemma and harmonic integration give
\begin{equation}\label{eq:depthlower}
 \Risk\ge\frac{c_*'}z
 \ln\left(1+\frac18\min\left\{\frac{z}{M\beta_0},\frac Mz\right\}\right),
 \qquad c_*'=\frac1{8192A_{\rm box}}.
\end{equation}
For $z>M/4$, the displayed right side is at most
$c_*'M/(8z^2)<2c_*'/M$, which is dominated by the full-history stationary IID lower $1/(128\sqrt{M+2})$. Thus \eqref{eq:depthlower} holds for every $z\ge1$; the large-$z$ case does not posit nonexistent rare windows. For the growing-depth finite evaluator above, take $E=1$.

Finally set $\beta=0$ and $c_-\sqrt M\le Q\le c_+\sqrt M$. If $\gamma\ge M^{1/4}$, then already $\gamma\gtrsim\ln M$. Otherwise,
\[
 z\ge c_-\sqrt M,\qquad M/z\ge M^{1/4}/c_+,
\]
so the logarithm in \eqref{eq:depthlower} is at least $(\ln M)/8$ for all sufficiently large $M$. A risk bound $C_{\rm risk}/\sqrt M$ forces
$\gamma\ge[c_*'/(8c_+C_{\rm risk})]\ln M$.
Together with \eqref{eq:depthgain}, this gives
\[
 d\ge\frac{\ln\ln M+O(1)}{\ln C_*}.
\]
The constants depend only on the fixed primitive budgets, band and risk multiplier. This proves the necessary half of Theorem~\ref{thm:depth}; the preceding construction proves sufficiency. The lower sources are stationary, hence also admissible in the arbitrary-start minimax problem.

\subsection{Paid depth and arithmetic ledgers}\label{sec:depthcost}
There are two complete implementation choices. They should not be combined coordinatewise into an artificial smaller budget.

For the generic growing-depth FA evaluator, let $U$ count supplied numerical records across \emph{all} stages and $F$ count all public mask/metadata bits; neither is depth-independent. With $B_p$ the maximum supplied description length and $n=M+3$, a guarded word can have
\[
 w=O_{\rm fixed}\bigl(P+\log(n+1)+\log(d+1)+\log(P+1)+\log(U+F+2)\bigr).
\]
Acquisition pays $UB_p+F$ bits, conversion pays $O(U(B_p+P)^2)$ work, and converted records occupy $O(UP)$ bits. Fixed arithmetic microcode has a safe $O(w^4)$ charge. Online sequential work is
\[
 O_{\rm fixed}\bigl(dn^2(P+1)w^2+(U+F)w+w^4\bigr),
\]
including all edges, Taylor steps, primitive updates, local input acquisition and tape rewinds. A conservative phase-sum peak is
\[
 O_{\rm fixed}\bigl(UB_p+UP+F+n+nw+w^4+B_p+P\bigr).
\]
Only numerical arrays roll in $O(nw)$ space; total parameter and mask storage still grows with the stages. Supplied-data generators add their actual work and scratch. The inverse/normalization decoder and breakpoint generator are additional costs, not already absorbed by this encoder ledger.

For the specialized attention-free compander stages, put
$\ell=\ceil{\log_2(n+K+2)}$ and $b=128(h+\ell+L+1)$.
Persistent scalars use $b$ bits, full products use $2b$, and guarded temporaries at most $4b+4$. A conservative fixed-loop budget is
\begin{equation}\label{eq:compandercost}
 \begin{aligned}
 T_0&=O\bigl([n^2K^6+n(L+1)K^6+(L+1)^2]b^2\bigr),\\
 S_0&=O\bigl(n^2+nK^3b+(L+1)K^6b+b^2\bigr).
 \end{aligned}
\end{equation}
These charges include dense zero entries, paid local features, all-token histograms and companders, complete output banks, decoder scans, parameter generation, masks, controllers, copies, clears and scratch. Stage parameters cost $O((L+1)K^6b)$ bits. If dummy dense heads are actually run at every stage, replace the first time term by $n^2(L+1)K^6$ and charge any duplicated mask storage.

The inverse is finite exact arithmetic: generate the $L+2$ endpoints $0,4^{-L},\ldots,4^{-1},1$, evaluate their dyadic $g_L$ values, scan every interval with a fixed tie rule, and solve its affine formula. The generation cost $O((L+1)^2b^2)$ is included in \eqref{eq:compandercost}. Inverse outputs have a common dyadic denominator dividing $2^{h+2L+O(1)}$. Normalization, positive fallback and final ordered rounding then use exact integers; no square-root, logarithm or exponential oracle is present.

An addressed evaluator stores $O(T_0b)$ instruction bits, uses
$O(T_0S_0b)$ forward work and $O(T_0b+S_0)$ peak, and pays
$O((T_0+S_0)b^2)$ offline work.
A separate PURE expansion with $T=T_0+b$, $S'=S_0+b$ has depth $O(T)$, width $O(S')$, sparse addressed coefficient bits $O(bTS')$, dense coefficient bits $O(bT(S')^2)$ and dense work $O(b^2T(S')^2)$. Phase peaks add the chosen program or coefficients, masks, input, output and rolling storage. At the statistical operating point $P=O(\log M)$ and $b=O(\log M+\log K)$: the gain in primitive depth does not make internal precision or total memory constant.

\section{Finite resource interpretation}\label{sec:costs}
The theorems distinguish information at the cut from the cost of producing it. None is a total-memory or runtime optimality claim. The fixed-depth and one-block ledgers below complement the genuinely depth-dependent ledgers in Section~\ref{sec:depthcost}. They illustrate what is actually paid, not an optimal implementation claim.

For the affine/ReLU evaluator, let $U$ be the actual number of supplied numerical records, $B_{\mathrm{param}}$ their maximum description length, and $F$ the explicit masks and other public-data bits. With $n=M+3$, a full signed online word has
\[
 w=O_{\mathrm{arch}}(P+\log n+\log(P+1)),
\]
augmented by $\log(U+F+2)$ when needed for addresses. Converted public parameter records occupy compact $P+O_{\mathrm{arch}}(1)$ slots, not $w$-bit workspace slots. Their total is $O_{\mathrm{arch}}(UP)$; padding them to $w$ bits is a separate format with $O(Uw)$ storage and writes. Original descriptions remain a separate $UB_{\mathrm{param}}$ charge.

One fixed-control multitape schedule scans position-ordered projected arrays, streams row-major masks and biases, and rewinds key/value tapes for every query row. All edges, including masked ones, are checked. Exact products, integer Taylor steps and division give online bitwork
\[
 O_{\mathrm{arch}}(n^2(P+1)w^2),
\]
plus supplied-parameter conversion $O(U(B_{\mathrm{param}}+P)^2)$. A safe retained-data and working-memory bound is
\[
 O_{\mathrm{arch}}(UB_{\mathrm{param}}+UP+F+n+nw+w^4+B_{\mathrm{param}}+P).
\]
These are phase sums including code, masks, input, both rolling arrays, output, controllers and scratch. A generator for supplied data adds its actual cost; arbitrary tables are not presumed free to generate. The compressed prototype also pays the finite log-constant generator, $O(p_0^3\log^2(p_0+1))$ work, and the full downstream decoder implementation. Shared buffers reduce peaks only under an explicit lifetime schedule.

For the actual triple-bank program, let $\ell=\ceil{\log_2(n+K+2)}$, $k=\bl(\kappa)$ and $b=64(h+\ell+k+1)$. Persistent words have $b$ bits; products retain their full $2b$ bits and guarded temporaries use at most $4b+4$. A conservative dense Boolean schedule has
\begin{equation}\label{eq:costtriple}
 T_0=O(n^2K^6b^2),\qquad
 S=O(n^2+nK^3b+K^6b+b^2).
\end{equation}
It pays the two singleton heads, all local gates, zero-inclusive dense maps, all-token $K^3$ banks, denominator and numerator accumulations, block selection, output normalization, copies, clears, masks, controllers and addresses. The cache consists of the actual integers $1,\kappa$. There is no numerical attention matrix hidden in the space bound.

An addressed sequential evaluator storing $O(T_0b)$ instruction bits can scan the $S$-bit tape for each instruction, costing $O(T_0Sb)$ forward bitwork and $O(T_0b+S)$ peak. Its offline generator pays $O((T_0+S)b^2)$ work. Alternatively, a pure Boolean/ReLU expansion with $T=T_0+b$, $S'=S+b$ has depth $O(T)$, width $O(S')$, $O(TS')$ gates and identity holds, $O(bT(S')^2)$ dense coefficient bits, and $O(b^2T(S')^2)$ dense work. The stored parameters, input, public constants, rolling layers and output must all be included in each phase peak. These are distinct complete implementations, not coordinates to combine into an artificial best-of-both estimate.

At $\kappa=\bl(M+1)$, the weight bitlength is $O(\log\log M)$, but $h$ and accumulation length are still $O(\log M)$. Effective score range $\ln\kappa$ is neither floating-point exponent precision nor the total bits used by the model. Likewise fixed ideal attention depth does not imply fixed depth after arithmetic unrolling.

\section{Proof map and limitations}\label{sec:map}
The manuscript is modular. It gives the main source construction, message argument, causal influence calculation, arithmetic mechanism, occupation transfer and score tradeoff, while relying on the explicitly identified statistical and implementation proof bodies. The accompanying dossier supplies full statements, proofs, dependency bindings, independent LLM reviews and finite diagnostics. Frozen source headers may describe an earlier candidate status; the accompanying acceptance records identify the adopted versions. An audit is evidence of checking, not formal verification or external peer review.

\begin{center}\small
\begin{longtable}{@{}>{\raggedright\arraybackslash}p{0.26\textwidth}>{\raggedright\arraybackslash}p{0.25\textwidth}>{\raggedright\arraybackslash}p{0.42\textwidth}@{}}
\toprule
Manuscript result & Dossier proof unit & Load-bearing scope\\\midrule\endhead
Theorem~\ref{thm:main}, original bank & \unit{FR001 v1}; \unit{SF001 v1}; \unit{AU001 v1} & Stationary all-$Q$ completion, arbitrary-start upper, subset/admissibility composition.\\[4pt]
Local-family lower and risk & \unit{BA001 v1}; \unit{BF001 v3} & Full bank plus boundary labels; positive bounded weights; actual finite-prefix inclusion and augmented costs.\\[4pt]
Causal influence and lower & \unit{CA001 v3}, boundary-causality scope note v1 & Matrix-entry envelope, square summability, fixed dense masks, formal frozen-branch domain; no future-dependent historical branches in the causal subclass.\\[4pt]
Theorem~\ref{thm:arithmetic} & \unit{FA001 v3} & Signed full products, positive-shift Taylor, bootstrap, sequential access and compact parameter slots.\\[4pt]
Theorem~\ref{thm:main}, finite causal family & \unit{CF001 v1} & Explicit four-head ideal prototype, compiled bank accuracy/range, downstream transfer and complete cost composition.\\[4pt]
Corollary~\ref{cor:threshold} & \unit{PF001 v1}; \unit{FR001 v1}; \unit{CF001 v1} & Uniform threshold from the all-resolution profile; known-promise budget selection only.\\[4pt]
Conservative finite decoder & \unit{NB002 v1}; \unit{SRC} & Exact master identities, $76/Q$ and $6/Q$ frequency errors, monotone rejection and event-radius interface.\\[4pt]
Equations \eqref{eq:inherited} and \eqref{eq:preint} & \unit{R011 v1}; \unit{R001}, \unit{S024}, \unit{S025}; \unit{NB002} & Inherited full-history statistics, ORIGINAL add-one and selector certificates; not reproved as a new statistical theorem.\\[4pt]
Terminal occupation & \unit{G/U} Section 11.1; \unit{IU001 v1}; \unit{AU001 v1} & Arbitrary-start first-occurrence counting; conditional return/singular-norm proof; complete residual suffix mixture.\\[4pt]
Theorem~\ref{thm:scoreupper} & \unit{SB001 v2}; \unit{GW001 v3} & Actual $K^3$ bank and integer weights; unchanged finite ledger and inherited ORIGINAL comparator.\\[4pt]
Proposition~\ref{prop:rescaled}; Theorem~\ref{thm:kappa} & \unit{SH001 v1}; \unit{SB001 v2} & Rescaled rare-window tube with a $\kappa$-independent constant; uniform real-budget infimum lower and integer upper in a fixed square-root grid band.\\[4pt]
Corollary~\ref{cor:hist} & \unit{SB001 v2} Section 7 at $\kappa=1$; \unit{BA001 v1} & Manuscript synthesis with the two proof lines shown; not a separate accepted unit.\\[4pt]
Proposition~\ref{prop:iiduniform} & \unit{UI001 v1} & Fixed heavy-window oracle, partial-reveal variance charge and posterior-density control; every finite weight envelope.\\[4pt]
Theorem~\ref{thm:joint}; Corollary~\ref{cor:jointthreshold} & \unit{TR001 v1} & UI001 and SH001 plus exhaustive joint-regime completion and unchanged SB001 upper; $K\ge3$, both initializations.\\[4pt]
Proposition~\ref{prop:binarylower} & \unit{BI001 v2} & Actual binary stationary witness, reference-path domination, Markov coupling and informative endpoint.\\[4pt]
Lemma~\ref{lem:binaryreturn}; Proposition~\ref{prop:binaryupper} & \unit{BU001 v2} & Stationary weighted returns; zero-weight support and zero-bank floor branches; unchanged eight-coordinate program.\\[4pt]
Theorem~\ref{thm:binaryjoint}; Corollary~\ref{cor:alphabet} & \unit{BS001 v1} & All-regime stationary binary synthesis, separate grid/product thresholds and fixed-architecture alphabet contrast.\\[4pt]
Proposition~\ref{prop:mixingfailure} & \unit{WM001 v2} & Full-support weighted endpoint obstruction; negative control for a transfer method, not a precision lower.\\[4pt]
Absorbing arbitrary starts & \unit{BAI002 v1}; \unit{BAI003 v2} & Restricted active-core starting point and unrestricted singleton-support run-PGF extension; true exit weights and original-path domination.\\[4pt]
Lemma~\ref{lem:schur}; diagonal/crossing bounds & \unit{BAI004 v5} & Schur reciprocal identity, initial-delay bridge, crossing renewal and weak-certificate/zero-exit branches.\\[4pt]
Equation~\eqref{eq:oppositediag} & \unit{BOD001 v1} & Opposite-diagonal first-hit density and weighted hitting potential, valid without a strong-certificate horizon restriction.\\[4pt]
Theorem~\ref{thm:binaryall} & \unit{BIF001 v1} & Complete support/initialization assembly and unchanged binary floor decoder; both-initialization thresholds.\\[4pt]
Theorem~\ref{thm:binarycausal} & \unit{BC001 v1}; \unit{BIF001 v1} & Markov nonlinear-message lower; exact finite histogram labels; distinct causal encoder and all-start composition.\\[4pt]
Theorem~\ref{thm:depth}, upper & \unit{GD001} upper derivation & Dyadic compander, downward purity, inverse normalization, fractional occupation and paid program.\\[4pt]
Theorem~\ref{thm:depth}, lower & \unit{GD001} lower derivation & Fixed per-stage budgets, frozen terminal initialization, depth-independent rescaled tube; $K\ge3$, $\beta=0$ necessity.\\[4pt]
Equation~\eqref{eq:depthprecision} & \unit{GD001} uniform FA extension & Same finite evaluator, uniform $P=4h+O(d+1)$ certificate, upper inclusion and full depth-dependent ledger.\\

\bottomrule
\end{longtable}
\end{center}

The source manifest shipped with the editable manuscript identifies exact paths and SHA-256 hashes for these proof versions. The earlier complete statistical dossier is necessary for the full dependency chain of \unit{R011}; this article does not relabel that inherited result as a new standalone proof. Its universal occupation argument is reproduced, but the all-parameter add-one and adaptive-selector statistical estimates remain inherited inputs.

Several boundaries are substantive. The fixed-family lower constants depend on architecture size. The growing-depth theorem removes dependence on the number of stages from the tube constant only by keeping every primitive, width, locality and boundary budget uniform. It does not allow arbitrary complexity inside one nominal layer. The one-block grid/weight frontier is complete in its declared variables, but no joint all-resolution grid/range/depth frontier is proved. The necessary depth order is established at $\beta=0$, $K\ge3$ and $Q\asymp\sqrt M$, not for every alphabet or positive mixing promise.

Query-selective hard masks, unbounded values, unrestricted nonlocal preprocessing, vanishing normalization constants and extra exact ports change the encoder class. Post-cut decoding is already unrestricted; extra computation only there cannot overcome the information restriction. A rounded program without the uniform ideal-to-finite certificate is not covered by the causal lower merely because some of its coefficients are bounded.

The fixed-$K$ rates do not establish sharp joint-alphabet behavior. Binary one-block and fixed-depth causal frontiers are now closed for both stationary and arbitrary initialization, including the square-root cutoff; the former stationary-only limitation is removed by Theorem~\ref{thm:binaryall}, not by endpoint mixing. The general-$K$ compander upper uses the inherited statistical input and is valid without a mixing promise, but supplies only a sufficient bound away from the matched depth operating point. Other losses, higher-order or nonhomogeneous sources, optimal total-resource tradeoffs and physically noisy computation are not settled.

The constructions are not claims about convergence of training, successful emitted-network execution or practical numerical constants. Companding, weighted counting and general depth--precision ideas are established mechanisms; this manuscript does not certify novelty. Direct counting with full-history access remains an important paid baseline, and no attention-specific computational advantage follows.

The resulting picture has three distinct resource coordinates. A finite grid loses rare-context mass; a larger pre-cut weight ratio can protect it within a one-block bank; repeated bounded nonlinear stages can instead redistribute grid resolution with fixed score range. Which protection is necessary depends on alphabet, initialization, mixing promise and the exact encoder contract. The modular proofs and finite ledgers make those dependencies explicit without turning them into universal architectural claims.

\begingroup
\small
\setlength{\bibsep}{4pt}
\bibliographystyle{plainnat}
\bibliography{references}
\endgroup
\end{document}
